\documentclass[
  doc,
  floatsintext,
  longtable,
  a4paper,
  nolmodern,
  notxfonts,
  notimes,
  12pt,
  colorlinks=true,linkcolor=blue,citecolor=blue,urlcolor=blue]{apa7}

\usepackage{amsmath}
\usepackage{amssymb}

\geometry{inner=1in, outer=1in}
\fancyhfoffset[LE,RO]{0cm}


\usepackage[bidi=default]{babel}
\babelprovide[main,import]{spanish}


% get rid of language-specific shorthands (see #6817):
\let\LanguageShortHands\languageshorthands
\def\languageshorthands#1{}

\RequirePackage{longtable}
\RequirePackage{threeparttablex}

\makeatletter
\renewcommand{\paragraph}{\@startsection{paragraph}{4}{\parindent}%
	{0\baselineskip \@plus 0.2ex \@minus 0.2ex}%
	{-.5em}%
	{\normalfont\normalsize\bfseries\typesectitle}}

\renewcommand{\subparagraph}[1]{\@startsection{subparagraph}{5}{0.5em}%
	{0\baselineskip \@plus 0.2ex \@minus 0.2ex}%
	{-\z@\relax}%
	{\normalfont\normalsize\bfseries\itshape\hspace{\parindent}{#1}\textit{\addperi}}{\relax}}
\makeatother




\usepackage{longtable, booktabs, multirow, multicol, colortbl, hhline, caption, array, float, xpatch}
\usepackage{subcaption}


\renewcommand\thesubfigure{\Alph{subfigure}}
\setcounter{topnumber}{2}
\setcounter{bottomnumber}{2}
\setcounter{totalnumber}{4}
\renewcommand{\topfraction}{0.85}
\renewcommand{\bottomfraction}{0.85}
\renewcommand{\textfraction}{0.15}
\renewcommand{\floatpagefraction}{0.7}

\usepackage{tcolorbox}
\tcbuselibrary{listings,theorems, breakable, skins}
\usepackage{fontawesome5}

\definecolor{quarto-callout-color}{HTML}{909090}
\definecolor{quarto-callout-note-color}{HTML}{0758E5}
\definecolor{quarto-callout-important-color}{HTML}{CC1914}
\definecolor{quarto-callout-warning-color}{HTML}{EB9113}
\definecolor{quarto-callout-tip-color}{HTML}{00A047}
\definecolor{quarto-callout-caution-color}{HTML}{FC5300}
\definecolor{quarto-callout-color-frame}{HTML}{ACACAC}
\definecolor{quarto-callout-note-color-frame}{HTML}{4582EC}
\definecolor{quarto-callout-important-color-frame}{HTML}{D9534F}
\definecolor{quarto-callout-warning-color-frame}{HTML}{F0AD4E}
\definecolor{quarto-callout-tip-color-frame}{HTML}{02B875}
\definecolor{quarto-callout-caution-color-frame}{HTML}{FD7E14}

%\newlength\Oldarrayrulewidth
%\newlength\Oldtabcolsep


\usepackage{hyperref}



\usepackage{color}
\usepackage{fancyvrb}
\newcommand{\VerbBar}{|}
\newcommand{\VERB}{\Verb[commandchars=\\\{\}]}
\DefineVerbatimEnvironment{Highlighting}{Verbatim}{commandchars=\\\{\}}
% Add ',fontsize=\small' for more characters per line
\usepackage{framed}
\definecolor{shadecolor}{RGB}{241,243,245}
\newenvironment{Shaded}{\begin{snugshade}}{\end{snugshade}}
\newcommand{\AlertTok}[1]{\textcolor[rgb]{0.68,0.00,0.00}{#1}}
\newcommand{\AnnotationTok}[1]{\textcolor[rgb]{0.37,0.37,0.37}{#1}}
\newcommand{\AttributeTok}[1]{\textcolor[rgb]{0.40,0.45,0.13}{#1}}
\newcommand{\BaseNTok}[1]{\textcolor[rgb]{0.68,0.00,0.00}{#1}}
\newcommand{\BuiltInTok}[1]{\textcolor[rgb]{0.00,0.23,0.31}{#1}}
\newcommand{\CharTok}[1]{\textcolor[rgb]{0.13,0.47,0.30}{#1}}
\newcommand{\CommentTok}[1]{\textcolor[rgb]{0.37,0.37,0.37}{#1}}
\newcommand{\CommentVarTok}[1]{\textcolor[rgb]{0.37,0.37,0.37}{\textit{#1}}}
\newcommand{\ConstantTok}[1]{\textcolor[rgb]{0.56,0.35,0.01}{#1}}
\newcommand{\ControlFlowTok}[1]{\textcolor[rgb]{0.00,0.23,0.31}{\textbf{#1}}}
\newcommand{\DataTypeTok}[1]{\textcolor[rgb]{0.68,0.00,0.00}{#1}}
\newcommand{\DecValTok}[1]{\textcolor[rgb]{0.68,0.00,0.00}{#1}}
\newcommand{\DocumentationTok}[1]{\textcolor[rgb]{0.37,0.37,0.37}{\textit{#1}}}
\newcommand{\ErrorTok}[1]{\textcolor[rgb]{0.68,0.00,0.00}{#1}}
\newcommand{\ExtensionTok}[1]{\textcolor[rgb]{0.00,0.23,0.31}{#1}}
\newcommand{\FloatTok}[1]{\textcolor[rgb]{0.68,0.00,0.00}{#1}}
\newcommand{\FunctionTok}[1]{\textcolor[rgb]{0.28,0.35,0.67}{#1}}
\newcommand{\ImportTok}[1]{\textcolor[rgb]{0.00,0.46,0.62}{#1}}
\newcommand{\InformationTok}[1]{\textcolor[rgb]{0.37,0.37,0.37}{#1}}
\newcommand{\KeywordTok}[1]{\textcolor[rgb]{0.00,0.23,0.31}{\textbf{#1}}}
\newcommand{\NormalTok}[1]{\textcolor[rgb]{0.00,0.23,0.31}{#1}}
\newcommand{\OperatorTok}[1]{\textcolor[rgb]{0.37,0.37,0.37}{#1}}
\newcommand{\OtherTok}[1]{\textcolor[rgb]{0.00,0.23,0.31}{#1}}
\newcommand{\PreprocessorTok}[1]{\textcolor[rgb]{0.68,0.00,0.00}{#1}}
\newcommand{\RegionMarkerTok}[1]{\textcolor[rgb]{0.00,0.23,0.31}{#1}}
\newcommand{\SpecialCharTok}[1]{\textcolor[rgb]{0.37,0.37,0.37}{#1}}
\newcommand{\SpecialStringTok}[1]{\textcolor[rgb]{0.13,0.47,0.30}{#1}}
\newcommand{\StringTok}[1]{\textcolor[rgb]{0.13,0.47,0.30}{#1}}
\newcommand{\VariableTok}[1]{\textcolor[rgb]{0.07,0.07,0.07}{#1}}
\newcommand{\VerbatimStringTok}[1]{\textcolor[rgb]{0.13,0.47,0.30}{#1}}
\newcommand{\WarningTok}[1]{\textcolor[rgb]{0.37,0.37,0.37}{\textit{#1}}}

\providecommand{\tightlist}{%
  \setlength{\itemsep}{0pt}\setlength{\parskip}{0pt}}
\usepackage{longtable,booktabs,array}
\usepackage{calc} % for calculating minipage widths
% Correct order of tables after \paragraph or \subparagraph
\usepackage{etoolbox}
\makeatletter
\patchcmd\longtable{\par}{\if@noskipsec\mbox{}\fi\par}{}{}
\makeatother
% Allow footnotes in longtable head/foot
\IfFileExists{footnotehyper.sty}{\usepackage{footnotehyper}}{\usepackage{footnote}}
\makesavenoteenv{longtable}

\usepackage{graphicx}
\makeatletter
\newsavebox\pandoc@box
\newcommand*\pandocbounded[1]{% scales image to fit in text height/width
  \sbox\pandoc@box{#1}%
  \Gscale@div\@tempa{\textheight}{\dimexpr\ht\pandoc@box+\dp\pandoc@box\relax}%
  \Gscale@div\@tempb{\linewidth}{\wd\pandoc@box}%
  \ifdim\@tempb\p@<\@tempa\p@\let\@tempa\@tempb\fi% select the smaller of both
  \ifdim\@tempa\p@<\p@\scalebox{\@tempa}{\usebox\pandoc@box}%
  \else\usebox{\pandoc@box}%
  \fi%
}
% Set default figure placement to htbp
\def\fps@figure{htbp}
\makeatother







\usepackage{newtx}

\defaultfontfeatures{Scale=MatchLowercase}
\defaultfontfeatures[\rmfamily]{Ligatures=TeX,Scale=1}





\title{Guía Completa de Positron IDE}


\shorttitle{Guía Completa de Positron IDE}


\usepackage{etoolbox}



\ccoppy{\textcopyright~2025}



\author{Edison Achalma}



\affiliation{
{Escuela Profesional de Economía, Universidad Nacional de San Cristóbal
de Huamanga}}




\leftheader{Achalma}

\date{2025-12-31}


\abstract{Positron es un IDE moderno y especializado en ciencia de
datos, construido sobre el núcleo open-source de VS Code (Code OSS) y
desarrollado por Posit Software. Ofrece soporte nativo para Python y R,
consola interactiva avanzada, paneles especializados (Variables, Data
Explorer, Plots, Help), integración completa con Quarto y herramientas
de asistencia IA. Esta guía completa cubre desde la instalación
multiplataforma hasta la configuración avanzada, migración desde VS
Code, personalización profunda y mejores prácticas para flujos de
trabajo intensivos en análisis de datos, programación estadística y
publicación científica reproducible. }

\keywords{Positron IDE, Ciencia de datos, Python, R, Quarto}

\authornote{\par{\addORCIDlink{Edison Achalma}{0000-0001-6996-3364}} 
\par{ }
\par{   El autor no tiene conflictos de interés que revelar.    Los
roles de autor se clasificaron utilizando la taxonomía de roles de
colaborador (CRediT; https://credit.niso.org/) de la siguiente
manera:  Edison Achalma:   conceptualización, metodología, análisis
formal, investigación, recursos, curación de
datos, redacción, visualización, supervisión, administración del
proyecto}
\par{La correspondencia relativa a este artículo debe dirigirse a Edison
Achalma, Escuela Profesional de Economía, Universidad Nacional de San
Cristóbal de Huamanga, Portal Independencia N°
57, Ayacucho, AYA 5001, Perú, Email: \href{mailto:elmer.achalma.09@unsch.edu.pe}{elmer.achalma.09@unsch.edu.pe}}
}

\makeatletter
\let\endoldlt\endlongtable
\def\endlongtable{
\hline
\endoldlt
}
\makeatother

\urlstyle{same}



\makeatletter
\@ifpackageloaded{caption}{}{\usepackage{caption}}
\AtBeginDocument{%
\ifdefined\contentsname
  \renewcommand*\contentsname{Tabla de contenidos}
\else
  \newcommand\contentsname{Tabla de contenidos}
\fi
\ifdefined\listfigurename
  \renewcommand*\listfigurename{Lista de Figuras}
\else
  \newcommand\listfigurename{Lista de Figuras}
\fi
\ifdefined\listtablename
  \renewcommand*\listtablename{Lista de Tablas}
\else
  \newcommand\listtablename{Lista de Tablas}
\fi
\ifdefined\figurename
  \renewcommand*\figurename{Figura}
\else
  \newcommand\figurename{Figura}
\fi
\ifdefined\tablename
  \renewcommand*\tablename{Tabla}
\else
  \newcommand\tablename{Tabla}
\fi
}
\@ifpackageloaded{float}{}{\usepackage{float}}
\floatstyle{ruled}
\@ifundefined{c@chapter}{\newfloat{codelisting}{h}{lop}}{\newfloat{codelisting}{h}{lop}[chapter]}
\floatname{codelisting}{Listado}
\newcommand*\listoflistings{\listof{codelisting}{Lista de Listados}}
\makeatother
\makeatletter
\makeatother
\makeatletter
\@ifpackageloaded{caption}{}{\usepackage{caption}}
\@ifpackageloaded{subcaption}{}{\usepackage{subcaption}}
\makeatother
\makeatletter
\@ifpackageloaded{fontawesome5}{}{\usepackage{fontawesome5}}
\makeatother

% From https://tex.stackexchange.com/a/645996/211326
%%% apa7 doesn't want to add appendix section titles in the toc
%%% let's make it do it
\makeatletter
\xpatchcmd{\appendix}
  {\par}
  {\addcontentsline{toc}{section}{\@currentlabelname}\par}
  {}{}
\makeatother

%% Disable longtable counter
%% https://tex.stackexchange.com/a/248395/211326

\usepackage{etoolbox}

\makeatletter
\patchcmd{\LT@caption}
  {\bgroup}
  {\bgroup\global\LTpatch@captiontrue}
  {}{}
\patchcmd{\longtable}
  {\par}
  {\par\global\LTpatch@captionfalse}
  {}{}
\apptocmd{\endlongtable}
  {\ifLTpatch@caption\else\addtocounter{table}{-1}\fi}
  {}{}
\newif\ifLTpatch@caption
\makeatother

\begin{document}

\maketitle


\hypertarget{toc}{}
\tableofcontents
\newpage
\section[Introduction]{Guía Completa de Positron IDE}

\setcounter{secnumdepth}{3}

\setlength\LTleft{0pt}




\pandocbounded{\includegraphics[keepaspectratio]{featured.png}}

\textbf{¿Qué es Positron?}

Positron es un IDE moderno y extensible para ciencia de datos construido
sobre \textbf{Code OSS} (el núcleo open-source de VS Code). Desarrollado
por \textbf{Posit Software} (anteriormente RStudio), está diseñado
específicamente para flujos de trabajo de ciencia de datos.

\textbf{Características Principales:}

\begin{itemize}
\tightlist
\item
  \textbf{Basado en Code OSS}: Toda la potencia de VS Code
\item
  \textbf{Soporte nativo Python y R}: Sin necesidad de extensiones
  externas
\item
  \textbf{Console interactiva}: IPython/IRKernel integrada
\item
  \textbf{Data Explorer}: Visualización de dataframes en vivo
\item
  \textbf{Variables Pane}: Inspección de variables
\item
  \textbf{Plots Pane}: Visualización de gráficos
\item
  \textbf{Help Pane}: Documentación integrada
\item
  \textbf{Connections Pane}: Gestión de bases de datos
\item
  \textbf{Quarto integrado}: Documentos dinámicos
\item
  \textbf{AI Assistant}: Positron Assistant y Databot
\item
  \textbf{Open VSX}: Marketplace de extensiones open-source
\end{itemize}

\textbf{Positron vs VS Code vs RStudio}

\begin{longtable}[]{@{}
  >{\raggedright\arraybackslash}p{(\linewidth - 6\tabcolsep) * \real{0.2895}}
  >{\raggedright\arraybackslash}p{(\linewidth - 6\tabcolsep) * \real{0.2368}}
  >{\raggedright\arraybackslash}p{(\linewidth - 6\tabcolsep) * \real{0.2368}}
  >{\raggedright\arraybackslash}p{(\linewidth - 6\tabcolsep) * \real{0.2368}}@{}}
\toprule\noalign{}
\begin{minipage}[b]{\linewidth}\raggedright
Característica
\end{minipage} & \begin{minipage}[b]{\linewidth}\raggedright
Positron
\end{minipage} & \begin{minipage}[b]{\linewidth}\raggedright
VS Code
\end{minipage} & \begin{minipage}[b]{\linewidth}\raggedright
RStudio
\end{minipage} \\
\midrule\noalign{}
\endhead
\bottomrule\noalign{}
\endlastfoot
Base técnica & Code OSS & Propietario + Code OSS & Propietario
(Posit) \\
Soporte nativo Python & Sí & No (extensión) & Limitado \\
Soporte nativo R & Sí & No (extensión) & Sí (muy optimizado) \\
Consola interactiva & Sí & No & Sí \\
Panel de datos/variables/plots & Sí & No & Sí \\
Soporte Quarto & Sí (nativo) & Sí (extensión) & Sí \\
Extensiones & Open VSX & VS Marketplace & Limitadas \\
Asistente IA integrado & Sí & No (requiere Copilot) & No \\
Multilenguaje & Sí (R + Python) & Sí (muy amplio) & Limitado (foco en
R) \\
Estabilidad & Alta & Muy alta & Muy alta \\
Crashes afectan el IDE & No & No & Sí \\
\end{longtable}

\begin{itemize}
\tightlist
\item
  \textbf{Positron} se posiciona como el sucesor natural para flujos de
  trabajo modernos de ciencia de datos que usan \textbf{R y Python} por
  igual. Combina lo mejor de RStudio (consola interactiva, panes de
  datos/plots/variables) con la extensibilidad y modernidad de VS Code.
\item
  \textbf{VS Code} sigue siendo el rey de la versatilidad general, pero
  requiere mucha configuración para obtener una experiencia fluida en R
  o Python data science.
\item
  \textbf{RStudio} permanece muy bien mantenido por Posit, con
  optimizaciones específicas para R que aún no están al 100\% en
  Positron (ej. algunos add-ins, soporte completo de .Rproj en ciertos
  casos, o features muy nicho como Sweave o ciertos publishing tools).
  Si tu trabajo es 95--100\% R y no necesitas Python, RStudio sigue
  siendo excelente y estable.
\end{itemize}

\section{Instalación}\label{instalaciuxf3n}

\subsection{Requisitos del Sistema}\label{requisitos-del-sistema}

\textbf{Mínimos:}

\begin{itemize}
\tightlist
\item
  Sistema Operativo: Linux, macOS, Windows
\item
  Espacio en disco: 500 MB
\item
  RAM: 2 GB (4 GB recomendado)
\end{itemize}

\subsection{Linux (Ubuntu/Debian)}\label{linux-ubuntudebian}

\begin{Shaded}
\begin{Highlighting}[]
\CommentTok{\# Descargar desde el sitio oficial}
\CommentTok{\# https://github.com/posit{-}dev/positron/releases}

\CommentTok{\# O instalar desde .deb}
\FunctionTok{sudo}\NormalTok{ dpkg }\AttributeTok{{-}i}\NormalTok{ positron\_}\PreprocessorTok{*}\NormalTok{.deb}
\FunctionTok{sudo}\NormalTok{ apt{-}get install }\AttributeTok{{-}f}  \CommentTok{\# Instalar dependencias}
\end{Highlighting}
\end{Shaded}

\subsection{macOS}\label{macos}

\begin{Shaded}
\begin{Highlighting}[]
\CommentTok{\# Descargar .dmg desde:}
\CommentTok{\# https://github.com/posit{-}dev/positron/releases}

\CommentTok{\# O con Homebrew}
\ExtensionTok{brew}\NormalTok{ install }\AttributeTok{{-}{-}cask}\NormalTok{ positron}
\end{Highlighting}
\end{Shaded}

\subsection{Windows}\label{windows}

\begin{Shaded}
\begin{Highlighting}[]
\CommentTok{\# Descargar instalador desde:}
\CommentTok{\# https://github.com/posit{-}dev/positron/releases}

\CommentTok{\# Ejecutar positron{-}setup.exe}
\end{Highlighting}
\end{Shaded}

\subsection{Verificar Instalación}\label{verificar-instalaciuxf3n}

\begin{Shaded}
\begin{Highlighting}[]
\CommentTok{\# Abrir Positron}
\ExtensionTok{positron}

\CommentTok{\# O desde línea de comandos}
\ExtensionTok{positron} \AttributeTok{{-}{-}version}
\end{Highlighting}
\end{Shaded}

\subsection{Configurar Comando en
Terminal}\label{configurar-comando-en-terminal}

\textbf{Linux/macOS:}

\begin{Shaded}
\begin{Highlighting}[]
\CommentTok{\# Añadir a \textasciitilde{}/.bashrc o \textasciitilde{}/.zshrc}
\BuiltInTok{export} \VariableTok{PATH}\OperatorTok{=}\StringTok{"}\VariableTok{$PATH}\StringTok{:/path/to/positron/bin"}

\CommentTok{\# O crear symlink}
\FunctionTok{sudo}\NormalTok{ ln }\AttributeTok{{-}s}\NormalTok{ /path/to/positron/bin/positron /usr/local/bin/positron}
\end{Highlighting}
\end{Shaded}

\section{Migración desde VS Code}\label{migraciuxf3n-desde-vs-code}

\subsection{Importar Configuración}\label{importar-configuraciuxf3n}

\textbf{Automáticamente al primer inicio:}

\begin{enumerate}
\def\labelenumi{\arabic{enumi}.}
\tightlist
\item
  Positron te preguntará si quieres importar settings de VS Code
\item
  Click en ``Yes'' para importar
\item
  Revisa los settings importados
\item
  Guarda los cambios
\end{enumerate}

\textbf{Manualmente:}

\begin{verbatim}
Command Palette (Ctrl+Shift+P) → Preferences: Import Settings...
\end{verbatim}

\subsection{Diferencias Clave}\label{diferencias-clave}

\textbf{Extensiones:}

\begin{itemize}
\tightlist
\item
  Positron usa \textbf{Open VSX} en lugar de VS Marketplace
\item
  Algunas extensiones pueden no estar disponibles
\item
  Python y R son nativos (no necesitas ms-python.python)
\end{itemize}

\textbf{Features Nuevas:}

\begin{itemize}
\tightlist
\item
  \textbf{Console}: Reemplaza el terminal para código interactivo
\item
  \textbf{Variables Pane}: Inspección de datos en vivo
\item
  \textbf{Data Explorer}: Visualización de dataframes
\item
  \textbf{Plots Pane}: Gráficos integrados
\item
  \textbf{Help Pane}: Documentación contextual
\end{itemize}

\subsection{Extensiones a Reemplazar}\label{extensiones-a-reemplazar}

\begin{longtable}[]{@{}ll@{}}
\toprule\noalign{}
VS Code & Positron \\
\midrule\noalign{}
\endhead
\bottomrule\noalign{}
\endlastfoot
ms-python.python & Built-in (nativo) \\
REditorSupport.r & Built-in (posit.air-vscode) \\
Jupyter & Built-in \\
GitLens & Instalar desde Open VSX \\
Prettier & Instalar desde Open VSX \\
\end{longtable}

\section{Configuración Inicial}\label{configuraciuxf3n-inicial}

\subsection{Ubicación de Archivos}\label{ubicaciuxf3n-de-archivos}

\begin{verbatim}
Linux: ~/.config/Positron/User/settings.json
macOS: ~/Library/Application Support/Positron/User/settings.json
Windows: %APPDATA%\Positron\User\settings.json
\end{verbatim}

\subsection{Abrir Settings}\label{abrir-settings}

\textbf{GUI:}

\begin{verbatim}
File → Preferences → Settings (Ctrl+,)
\end{verbatim}

\textbf{JSON:}

\begin{verbatim}
Command Palette → Preferences: Open User Settings (JSON)
\end{verbatim}

\subsection{Configuración Básica
Recomendada}\label{configuraciuxf3n-buxe1sica-recomendada}

\begin{Shaded}
\begin{Highlighting}[]
\FunctionTok{\{}
    \ErrorTok{//} \ErrorTok{Editor}
    \DataTypeTok{"editor.fontSize"}\FunctionTok{:} \DecValTok{14}\FunctionTok{,}
    \DataTypeTok{"editor.fontFamily"}\FunctionTok{:} \StringTok{"\textquotesingle{}JetBrains Mono\textquotesingle{}, \textquotesingle{}Fira Code\textquotesingle{}, \textquotesingle{}Hack Nerd Font\textquotesingle{}, monospace"}\FunctionTok{,}
    \DataTypeTok{"editor.fontLigatures"}\FunctionTok{:} \KeywordTok{true}\FunctionTok{,}
    \DataTypeTok{"editor.lineNumbers"}\FunctionTok{:} \StringTok{"on"}\FunctionTok{,}
    \DataTypeTok{"editor.minimap.enabled"}\FunctionTok{:} \KeywordTok{false}\FunctionTok{,}
    \DataTypeTok{"editor.formatOnSave"}\FunctionTok{:} \KeywordTok{true}\FunctionTok{,}
    \DataTypeTok{"editor.formatOnPaste"}\FunctionTok{:} \KeywordTok{true}\FunctionTok{,}
    \DataTypeTok{"editor.cursorBlinking"}\FunctionTok{:} \StringTok{"smooth"}\FunctionTok{,}
    \DataTypeTok{"editor.cursorSmoothCaretAnimation"}\FunctionTok{:} \StringTok{"on"}\FunctionTok{,}
    
    \ErrorTok{//} \ErrorTok{Tema}
    \DataTypeTok{"workbench.colorTheme"}\FunctionTok{:} \StringTok{"Default Positron Dark"}\FunctionTok{,}
    \DataTypeTok{"workbench.iconTheme"}\FunctionTok{:} \StringTok{"material{-}icon{-}theme"}\FunctionTok{,}
    
    \ErrorTok{//} \ErrorTok{Terminal}
    \DataTypeTok{"terminal.integrated.fontFamily"}\FunctionTok{:} \StringTok{"\textquotesingle{}Hack Nerd Font\textquotesingle{}, monospace"}\FunctionTok{,}
    \DataTypeTok{"terminal.integrated.fontSize"}\FunctionTok{:} \DecValTok{13}\FunctionTok{,}
    \DataTypeTok{"terminal.integrated.defaultProfile.linux"}\FunctionTok{:} \StringTok{"bash"}\FunctionTok{,}
    \DataTypeTok{"terminal.integrated.copyOnSelection"}\FunctionTok{:} \KeywordTok{true}\FunctionTok{,}
    \DataTypeTok{"terminal.integrated.cursorBlinking"}\FunctionTok{:} \KeywordTok{true}\FunctionTok{,}
    \DataTypeTok{"terminal.integrated.inheritEnv"}\FunctionTok{:} \KeywordTok{false}\FunctionTok{,}
    
    \ErrorTok{//} \ErrorTok{Git}
    \DataTypeTok{"git.enableSmartCommit"}\FunctionTok{:} \KeywordTok{true}\FunctionTok{,}
    \DataTypeTok{"git.autofetch"}\FunctionTok{:} \KeywordTok{true}\FunctionTok{,}
    \DataTypeTok{"git.confirmSync"}\FunctionTok{:} \KeywordTok{false}\FunctionTok{,}
    
    \ErrorTok{//} \ErrorTok{Files}
    \DataTypeTok{"files.autoSave"}\FunctionTok{:} \StringTok{"onFocusChange"}\FunctionTok{,}
    \DataTypeTok{"files.associations"}\FunctionTok{:} \FunctionTok{\{}
        \DataTypeTok{"*.qmd"}\FunctionTok{:} \StringTok{"quarto"}\FunctionTok{,}
        \DataTypeTok{"*.Rmd"}\FunctionTok{:} \StringTok{"rmarkdown"}\FunctionTok{,}
        \DataTypeTok{"renv.lock"}\FunctionTok{:} \StringTok{"json"}
    \FunctionTok{\},}
    
    \ErrorTok{//} \ErrorTok{Quarto}
    \DataTypeTok{"quarto.visualEditor.spellingDictionary"}\FunctionTok{:} \StringTok{"es\_ES"}\FunctionTok{,}
    
    \ErrorTok{//} \ErrorTok{Positron{-}specific}
    \DataTypeTok{"positron.assistant.enable"}\FunctionTok{:} \KeywordTok{true}\FunctionTok{,}
    \DataTypeTok{"positron.assistant.showTokenUsage.enable"}\FunctionTok{:} \KeywordTok{true}
\FunctionTok{\}}
\end{Highlighting}
\end{Shaded}

\section{Interfaz de Usuario}\label{interfaz-de-usuario}

\subsection{Layout Principal}\label{layout-principal}

\begin{verbatim}
┌─────────────────────────────────────────────────────┐
│ Title Bar                                           │
├──────────────────┬──────────────────────────────────┤
│                  │                                  │
│  Activity Bar    │                                  │
│  (left/bottom)   │       Editor Area                │
│                  │                                  │
│  - Explorer      │                                  │
│  - Search        ├──────────────────────────────────┤
│  - Source Ctrl   │                                  │
│  - Run/Debug     │       Console / Terminal         │
│  - Extensions    │                                  │
│                  │                                  │
├──────────────────┴──────────────────────────────────┤
│  Status Bar                                         │
└─────────────────────────────────────────────────────┘
\end{verbatim}

\subsection{Panes Exclusivos de
Positron}\label{panes-exclusivos-de-positron}

\textbf{Variables Pane:}

\begin{verbatim}
Muestra todas las variables en el entorno actual
- Nombre
- Tipo
- Valor/Tamaño
- Preview
\end{verbatim}

\textbf{Data Explorer:}

\begin{verbatim}
Visualización interactiva de dataframes
- Filtrado
- Ordenamiento
- Búsqueda
- Exportación
\end{verbatim}

\textbf{Plots Pane:}

\begin{verbatim}
Visualización de gráficos generados
- Historial de plots
- Zoom
- Exportar
\end{verbatim}

\textbf{Help Pane:}

\begin{verbatim}
Documentación contextual
- Búsqueda integrada
- Ejemplos
- Argumentos de funciones
\end{verbatim}

\textbf{Console:}

\begin{verbatim}
Shell interactivo para Python/R
- Completado de código
- Syntax highlighting
- Integración con Variables/Plots
\end{verbatim}

\subsection{Personalizar Layout}\label{personalizar-layout}

\textbf{Activity Bar Position:}

\begin{Shaded}
\begin{Highlighting}[]
\ErrorTok{"workbench.activityBar.location":} \ErrorTok{"bottom"} \ErrorTok{//} \ErrorTok{o} \ErrorTok{"top",} \ErrorTok{"left"}
\end{Highlighting}
\end{Shaded}

\textbf{Side Bar Position:}

\begin{Shaded}
\begin{Highlighting}[]
\ErrorTok{"workbench.sideBar.location":} \ErrorTok{"right"} \ErrorTok{//} \ErrorTok{o} \ErrorTok{"left"}
\end{Highlighting}
\end{Shaded}

\textbf{Command Palette Position:}

\begin{verbatim}
Arrastrarlo a la posición deseada
\end{verbatim}

\section{Atajos de Teclado}\label{atajos-de-teclado}

\subsection{Atajos Específicos de
Positron}\label{atajos-especuxedficos-de-positron}

\begin{longtable}[]{@{}ll@{}}
\toprule\noalign{}
Atajo & Acción \\
\midrule\noalign{}
\endhead
\bottomrule\noalign{}
\endlastfoot
\texttt{Ctrl+Enter} & Ejecutar línea/selección en Console \\
\texttt{Ctrl+Alt+Home} & Ejecutar desde inicio hasta línea actual \\
\texttt{Ctrl+Alt+End} & Ejecutar desde línea actual hasta fin \\
\texttt{Ctrl+Shift+0} & Reiniciar intérprete \\
\texttt{Ctrl+Shift+P} & Ejecutar archivo completo en Console \\
\texttt{F1} & Ayuda contextual \\
\texttt{Ctrl+K,\ Ctrl+R} & Ayuda contextual (alternativo) \\
\texttt{Ctrl+K,\ F} & Enfocar Console \\
\texttt{Ctrl+K,\ V} & Enfocar Variables pane \\
\texttt{Ctrl+L} & Limpiar Console \\
\end{longtable}

\subsection{Atajos Generales (heredados de VS
Code)}\label{atajos-generales-heredados-de-vs-code}

\textbf{Editor:}

\begin{longtable}[]{@{}ll@{}}
\toprule\noalign{}
Atajo & Acción \\
\midrule\noalign{}
\endhead
\bottomrule\noalign{}
\endlastfoot
\texttt{Ctrl+P} & Quick Open (buscar archivo) \\
\texttt{Ctrl+Shift+P} & Command Palette \\
\texttt{Ctrl+,} & Settings \\
\texttt{Ctrl+B} & Toggle Sidebar \\
\texttt{Ctrl+J} & Toggle Panel \\
\texttt{Ctrl+Shift+E} & Explorer \\
\texttt{Ctrl+Shift+F} & Search \\
\texttt{Ctrl+Shift+G} & Source Control \\
\texttt{Ctrl+Shift+D} & Run \& Debug \\
\texttt{Ctrl+Shift+X} & Extensions \\
\end{longtable}

\textbf{Edición:}

\begin{longtable}[]{@{}ll@{}}
\toprule\noalign{}
Atajo & Acción \\
\midrule\noalign{}
\endhead
\bottomrule\noalign{}
\endlastfoot
\texttt{Ctrl+Space} & Trigger suggestions \\
\texttt{Ctrl+Shift+K} & Delete line \\
\texttt{Alt+Up/Down} & Move line up/down \\
\texttt{Shift+Alt+Up/Down} & Copy line up/down \\
\texttt{Ctrl+/} & Toggle comment \\
\texttt{Ctrl+Shift+{[}} & Fold region \\
\texttt{Ctrl+Shift+{]}} & Unfold region \\
\texttt{Ctrl+K\ Ctrl+0} & Fold all \\
\texttt{Ctrl+K\ Ctrl+J} & Unfold all \\
\end{longtable}

\textbf{Navegación:}

\begin{longtable}[]{@{}ll@{}}
\toprule\noalign{}
Atajo & Acción \\
\midrule\noalign{}
\endhead
\bottomrule\noalign{}
\endlastfoot
\texttt{Ctrl+G} & Go to line \\
\texttt{Ctrl+Shift+O} & Go to symbol \\
\texttt{Ctrl+T} & Show all symbols \\
\texttt{F12} & Go to definition \\
\texttt{Alt+F12} & Peek definition \\
\texttt{Shift+F12} & Show references \\
\texttt{Ctrl+K\ F12} & Open definition to side \\
\end{longtable}

\textbf{Terminal:}

\begin{longtable}[]{@{}ll@{}}
\toprule\noalign{}
Atajo & Acción \\
\midrule\noalign{}
\endhead
\bottomrule\noalign{}
\endlastfoot
\texttt{Ctrl+\textasciigrave{}} & Toggle terminal \\
\texttt{Ctrl+Shift+\textasciigrave{}} & New terminal \\
\texttt{Ctrl+Shift+C} & Copy \\
\texttt{Ctrl+Shift+V} & Paste \\
\end{longtable}

\subsection{Personalizar Atajos}\label{personalizar-atajos}

\textbf{Abrir configuración:}

\begin{verbatim}
Command Palette → Preferences: Open Keyboard Shortcuts
\end{verbatim}

\textbf{Archivo JSON:}

\begin{verbatim}
Command Palette → Preferences: Open Keyboard Shortcuts (JSON)
\end{verbatim}

\textbf{Ejemplo de customización:}

\begin{Shaded}
\begin{Highlighting}[]
\OtherTok{[}
    \FunctionTok{\{}
        \DataTypeTok{"key"}\FunctionTok{:} \StringTok{"ctrl+shift+r"}\FunctionTok{,}
        \DataTypeTok{"command"}\FunctionTok{:} \StringTok{"workbench.action.executeCode.console"}\FunctionTok{,}
        \DataTypeTok{"when"}\FunctionTok{:} \StringTok{"editorTextFocus"}\FunctionTok{,}
        \DataTypeTok{"args"}\FunctionTok{:} \FunctionTok{\{}
            \DataTypeTok{"langId"}\FunctionTok{:} \StringTok{"r"}\FunctionTok{,}
            \DataTypeTok{"code"}\FunctionTok{:} \StringTok{"reprex::reprex\_selection()"}\FunctionTok{,}
            \DataTypeTok{"focus"}\FunctionTok{:} \KeywordTok{true}
        \FunctionTok{\}}
    \FunctionTok{\}}\OtherTok{,}
    \FunctionTok{\{}
        \DataTypeTok{"key"}\FunctionTok{:} \StringTok{"ctrl+alt+f"}\FunctionTok{,}
        \DataTypeTok{"command"}\FunctionTok{:} \StringTok{"workbench.action.formatDocument"}
    \FunctionTok{\}}\OtherTok{,}
    \FunctionTok{\{}
        \DataTypeTok{"key"}\FunctionTok{:} \StringTok{"f2"}\FunctionTok{,}
        \DataTypeTok{"command"}\FunctionTok{:} \StringTok{"positron.dataExplorer.openDataExplorer"}
    \FunctionTok{\}}
\OtherTok{]}
\end{Highlighting}
\end{Shaded}

\section{Configuración Avanzada}\label{configuraciuxf3n-avanzada}

\subsection{Settings.json Completo
(Base)}\label{settings.json-completo-base}

\begin{Shaded}
\begin{Highlighting}[]
\FunctionTok{\{}
    \ErrorTok{//} \ErrorTok{==========================================}
    \ErrorTok{//} \ErrorTok{EDITOR}
    \ErrorTok{//} \ErrorTok{==========================================}
    \DataTypeTok{"editor.fontSize"}\FunctionTok{:} \DecValTok{13}\FunctionTok{,}
    \DataTypeTok{"editor.fontFamily"}\FunctionTok{:} \StringTok{"\textquotesingle{}JetBrains Mono\textquotesingle{}, \textquotesingle{}Fira Code\textquotesingle{}, \textquotesingle{}Hack Nerd Font\textquotesingle{}, monospace"}\FunctionTok{,}
    \DataTypeTok{"editor.fontLigatures"}\FunctionTok{:} \KeywordTok{true}\FunctionTok{,}
    \DataTypeTok{"editor.lineNumbers"}\FunctionTok{:} \StringTok{"on"}\FunctionTok{,}
    \DataTypeTok{"editor.rulers"}\FunctionTok{:} \OtherTok{[}\DecValTok{80}\OtherTok{,} \DecValTok{120}\OtherTok{]}\FunctionTok{,}
    \DataTypeTok{"editor.wordWrap"}\FunctionTok{:} \StringTok{"off"}\FunctionTok{,}
    \DataTypeTok{"editor.minimap.enabled"}\FunctionTok{:} \KeywordTok{false}\FunctionTok{,}
    \DataTypeTok{"editor.scrollbar.vertical"}\FunctionTok{:} \StringTok{"auto"}\FunctionTok{,}
    \DataTypeTok{"editor.scrollbar.horizontal"}\FunctionTok{:} \StringTok{"hidden"}\FunctionTok{,}
    \DataTypeTok{"editor.cursorBlinking"}\FunctionTok{:} \StringTok{"smooth"}\FunctionTok{,}
    \DataTypeTok{"editor.cursorSmoothCaretAnimation"}\FunctionTok{:} \StringTok{"on"}\FunctionTok{,}
    \DataTypeTok{"editor.cursorStyle"}\FunctionTok{:} \StringTok{"line"}\FunctionTok{,}
    \DataTypeTok{"editor.formatOnSave"}\FunctionTok{:} \KeywordTok{true}\FunctionTok{,}
    \DataTypeTok{"editor.formatOnPaste"}\FunctionTok{:} \KeywordTok{true}\FunctionTok{,}
    \DataTypeTok{"editor.formatOnType"}\FunctionTok{:} \KeywordTok{false}\FunctionTok{,}
    \DataTypeTok{"editor.tabSize"}\FunctionTok{:} \DecValTok{4}\FunctionTok{,}
    \DataTypeTok{"editor.insertSpaces"}\FunctionTok{:} \KeywordTok{true}\FunctionTok{,}
    \DataTypeTok{"editor.detectIndentation"}\FunctionTok{:} \KeywordTok{true}\FunctionTok{,}
    \DataTypeTok{"editor.renderWhitespace"}\FunctionTok{:} \StringTok{"selection"}\FunctionTok{,}
    \DataTypeTok{"editor.bracketPairColorization.enabled"}\FunctionTok{:} \KeywordTok{true}\FunctionTok{,}
    \DataTypeTok{"editor.guides.bracketPairs"}\FunctionTok{:} \KeywordTok{true}\FunctionTok{,}
    \DataTypeTok{"editor.linkedEditing"}\FunctionTok{:} \KeywordTok{true}\FunctionTok{,}
    \DataTypeTok{"editor.suggest.preview"}\FunctionTok{:} \KeywordTok{true}\FunctionTok{,}
    \DataTypeTok{"editor.inlineSuggest.enabled"}\FunctionTok{:} \KeywordTok{true}\FunctionTok{,}
    \DataTypeTok{"editor.quickSuggestions"}\FunctionTok{:} \FunctionTok{\{}
        \DataTypeTok{"other"}\FunctionTok{:} \KeywordTok{true}\FunctionTok{,}
        \DataTypeTok{"comments"}\FunctionTok{:} \KeywordTok{false}\FunctionTok{,}
        \DataTypeTok{"strings"}\FunctionTok{:} \KeywordTok{false}
    \FunctionTok{\},}
    
    \ErrorTok{//} \ErrorTok{==========================================}
    \ErrorTok{//} \ErrorTok{WORKBENCH}
    \ErrorTok{//} \ErrorTok{==========================================}
    \DataTypeTok{"workbench.colorTheme"}\FunctionTok{:} \StringTok{"Default Positron Dark"}\FunctionTok{,}
    \DataTypeTok{"workbench.iconTheme"}\FunctionTok{:} \StringTok{"material{-}icon{-}theme"}\FunctionTok{,}
    \DataTypeTok{"workbench.productIconTheme"}\FunctionTok{:} \StringTok{"Default"}\FunctionTok{,}
    \DataTypeTok{"workbench.activityBar.location"}\FunctionTok{:} \StringTok{"bottom"}\FunctionTok{,}
    \DataTypeTok{"workbench.sideBar.location"}\FunctionTok{:} \StringTok{"left"}\FunctionTok{,}
    \DataTypeTok{"workbench.panel.defaultLocation"}\FunctionTok{:} \StringTok{"bottom"}\FunctionTok{,}
    \DataTypeTok{"workbench.statusBar.visible"}\FunctionTok{:} \KeywordTok{true}\FunctionTok{,}
    \DataTypeTok{"workbench.editor.showTabs"}\FunctionTok{:} \StringTok{"multiple"}\FunctionTok{,}
    \DataTypeTok{"workbench.editor.tabSizing"}\FunctionTok{:} \StringTok{"fit"}\FunctionTok{,}
    \DataTypeTok{"workbench.editor.highlightModifiedTabs"}\FunctionTok{:} \KeywordTok{true}\FunctionTok{,}
    \DataTypeTok{"workbench.editor.labelFormat"}\FunctionTok{:} \StringTok{"short"}\FunctionTok{,}
    \DataTypeTok{"workbench.editor.limit.enabled"}\FunctionTok{:} \KeywordTok{true}\FunctionTok{,}
    \DataTypeTok{"workbench.editor.limit.value"}\FunctionTok{:} \DecValTok{10}\FunctionTok{,}
    \DataTypeTok{"workbench.editor.empty.hint"}\FunctionTok{:} \StringTok{"hidden"}\FunctionTok{,}
    \DataTypeTok{"workbench.startupEditor"}\FunctionTok{:} \StringTok{"none"}\FunctionTok{,}
    
    \ErrorTok{//} \ErrorTok{==========================================}
    \ErrorTok{//} \ErrorTok{TERMINAL}
    \ErrorTok{//} \ErrorTok{==========================================}
    \DataTypeTok{"terminal.integrated.fontFamily"}\FunctionTok{:} \StringTok{"\textquotesingle{}Hack Nerd Font\textquotesingle{}, \textquotesingle{}FiraCode Nerd Font\textquotesingle{}, monospace"}\FunctionTok{,}
    \DataTypeTok{"terminal.integrated.fontSize"}\FunctionTok{:} \DecValTok{12}\FunctionTok{,}
    \DataTypeTok{"terminal.integrated.fontWeight"}\FunctionTok{:} \StringTok{"normal"}\FunctionTok{,}
    \DataTypeTok{"terminal.integrated.lineHeight"}\FunctionTok{:} \FloatTok{1.2}\FunctionTok{,}
    \DataTypeTok{"terminal.integrated.cursorBlinking"}\FunctionTok{:} \KeywordTok{true}\FunctionTok{,}
    \DataTypeTok{"terminal.integrated.cursorStyle"}\FunctionTok{:} \StringTok{"line"}\FunctionTok{,}
    \DataTypeTok{"terminal.integrated.defaultProfile.linux"}\FunctionTok{:} \StringTok{"bash"}\FunctionTok{,}
    \DataTypeTok{"terminal.integrated.copyOnSelection"}\FunctionTok{:} \KeywordTok{true}\FunctionTok{,}
    \DataTypeTok{"terminal.integrated.inheritEnv"}\FunctionTok{:} \KeywordTok{false}\FunctionTok{,}
    \DataTypeTok{"terminal.integrated.enablePersistentSessions"}\FunctionTok{:} \KeywordTok{false}\FunctionTok{,}
    \DataTypeTok{"terminal.integrated.scrollback"}\FunctionTok{:} \DecValTok{10000}\FunctionTok{,}
    
    \ErrorTok{//} \ErrorTok{==========================================}
    \ErrorTok{//} \ErrorTok{FILES}
    \ErrorTok{//} \ErrorTok{==========================================}
    \DataTypeTok{"files.autoSave"}\FunctionTok{:} \StringTok{"onFocusChange"}\FunctionTok{,}
    \DataTypeTok{"files.trimTrailingWhitespace"}\FunctionTok{:} \KeywordTok{true}\FunctionTok{,}
    \DataTypeTok{"files.insertFinalNewline"}\FunctionTok{:} \KeywordTok{true}\FunctionTok{,}
    \DataTypeTok{"files.trimFinalNewlines"}\FunctionTok{:} \KeywordTok{true}\FunctionTok{,}
    \DataTypeTok{"files.associations"}\FunctionTok{:} \FunctionTok{\{}
        \DataTypeTok{"*.qmd"}\FunctionTok{:} \StringTok{"quarto"}\FunctionTok{,}
        \DataTypeTok{"*.Rmd"}\FunctionTok{:} \StringTok{"rmarkdown"}\FunctionTok{,}
        \DataTypeTok{"*.renv.lock"}\FunctionTok{:} \StringTok{"json"}\FunctionTok{,}
        \DataTypeTok{"renv.lock"}\FunctionTok{:} \StringTok{"json"}\FunctionTok{,}
        \DataTypeTok{"*.ipynb"}\FunctionTok{:} \StringTok{"jupyter{-}notebook"}
    \FunctionTok{\},}
    \DataTypeTok{"files.exclude"}\FunctionTok{:} \FunctionTok{\{}
        \DataTypeTok{"**/.git"}\FunctionTok{:} \KeywordTok{true}\FunctionTok{,}
        \DataTypeTok{"**/.DS\_Store"}\FunctionTok{:} \KeywordTok{true}\FunctionTok{,}
        \DataTypeTok{"**/\_\_pycache\_\_"}\FunctionTok{:} \KeywordTok{true}\FunctionTok{,}
        \DataTypeTok{"**/.Rproj.user"}\FunctionTok{:} \KeywordTok{true}\FunctionTok{,}
        \DataTypeTok{"**/.ipynb\_checkpoints"}\FunctionTok{:} \KeywordTok{true}\FunctionTok{,}
        \DataTypeTok{"**/node\_modules"}\FunctionTok{:} \KeywordTok{true}
    \FunctionTok{\},}
    
    \ErrorTok{//} \ErrorTok{==========================================}
    \ErrorTok{//} \ErrorTok{GIT}
    \ErrorTok{//} \ErrorTok{==========================================}
    \DataTypeTok{"git.enableSmartCommit"}\FunctionTok{:} \KeywordTok{true}\FunctionTok{,}
    \DataTypeTok{"git.autofetch"}\FunctionTok{:} \KeywordTok{true}\FunctionTok{,}
    \DataTypeTok{"git.confirmSync"}\FunctionTok{:} \KeywordTok{false}\FunctionTok{,}
    \DataTypeTok{"git.openRepositoryInParentFolders"}\FunctionTok{:} \StringTok{"always"}\FunctionTok{,}
    \DataTypeTok{"git.suggestSmartCommit"}\FunctionTok{:} \KeywordTok{false}\FunctionTok{,}
    
    \ErrorTok{//} \ErrorTok{==========================================}
    \ErrorTok{//} \ErrorTok{SEARCH}
    \ErrorTok{//} \ErrorTok{==========================================}
    \DataTypeTok{"search.exclude"}\FunctionTok{:} \FunctionTok{\{}
        \DataTypeTok{"**/node\_modules"}\FunctionTok{:} \KeywordTok{true}\FunctionTok{,}
        \DataTypeTok{"**/bower\_components"}\FunctionTok{:} \KeywordTok{true}\FunctionTok{,}
        \DataTypeTok{"**/.git"}\FunctionTok{:} \KeywordTok{true}\FunctionTok{,}
        \DataTypeTok{"**/\_\_pycache\_\_"}\FunctionTok{:} \KeywordTok{true}\FunctionTok{,}
        \DataTypeTok{"**/.Rproj.user"}\FunctionTok{:} \KeywordTok{true}\FunctionTok{,}
        \DataTypeTok{"**/.ipynb\_checkpoints"}\FunctionTok{:} \KeywordTok{true}
    \FunctionTok{\},}
    
    \ErrorTok{//} \ErrorTok{==========================================}
    \ErrorTok{//} \ErrorTok{DIFF} \ErrorTok{EDITOR}
    \ErrorTok{//} \ErrorTok{==========================================}
    \DataTypeTok{"diffEditor.ignoreTrimWhitespace"}\FunctionTok{:} \KeywordTok{false}\FunctionTok{,}
    \DataTypeTok{"diffEditor.renderSideBySide"}\FunctionTok{:} \KeywordTok{true}\FunctionTok{,}
    \DataTypeTok{"diffEditor.wordWrap"}\FunctionTok{:} \StringTok{"off"}\FunctionTok{,}
    \DataTypeTok{"diffEditor.maxComputationTime"}\FunctionTok{:} \DecValTok{0}\FunctionTok{,}
    
    \ErrorTok{//} \ErrorTok{==========================================}
    \ErrorTok{//} \ErrorTok{LANGUAGE{-}SPECIFIC}
    \ErrorTok{//} \ErrorTok{==========================================}
    \DataTypeTok{"[python]"}\FunctionTok{:} \FunctionTok{\{}
        \DataTypeTok{"editor.defaultFormatter"}\FunctionTok{:} \StringTok{"charliermarsh.ruff"}\FunctionTok{,}
        \DataTypeTok{"editor.formatOnSave"}\FunctionTok{:} \KeywordTok{true}\FunctionTok{,}
        \DataTypeTok{"editor.codeActionsOnSave"}\FunctionTok{:} \FunctionTok{\{}
            \DataTypeTok{"source.organizeImports"}\FunctionTok{:} \StringTok{"explicit"}
        \FunctionTok{\},}
        \DataTypeTok{"diffEditor.ignoreTrimWhitespace"}\FunctionTok{:} \KeywordTok{false}
    \FunctionTok{\},}
    \DataTypeTok{"[r]"}\FunctionTok{:} \FunctionTok{\{}
        \DataTypeTok{"editor.tabSize"}\FunctionTok{:} \DecValTok{2}\FunctionTok{,}
        \DataTypeTok{"editor.insertSpaces"}\FunctionTok{:} \KeywordTok{true}
    \FunctionTok{\},}
    \DataTypeTok{"[quarto]"}\FunctionTok{:} \FunctionTok{\{}
        \DataTypeTok{"editor.defaultFormatter"}\FunctionTok{:} \StringTok{"quarto.quarto"}\FunctionTok{,}
        \DataTypeTok{"editor.wordWrap"}\FunctionTok{:} \StringTok{"on"}
    \FunctionTok{\},}
    \DataTypeTok{"[markdown]"}\FunctionTok{:} \FunctionTok{\{}
        \DataTypeTok{"editor.defaultFormatter"}\FunctionTok{:} \StringTok{"esbenp.prettier{-}vscode"}\FunctionTok{,}
        \DataTypeTok{"editor.wordWrap"}\FunctionTok{:} \StringTok{"on"}
    \FunctionTok{\},}
    \DataTypeTok{"[yaml]"}\FunctionTok{:} \FunctionTok{\{}
        \DataTypeTok{"editor.defaultFormatter"}\FunctionTok{:} \StringTok{"quarto.quarto"}\FunctionTok{,}
        \DataTypeTok{"editor.tabSize"}\FunctionTok{:} \DecValTok{2}
    \FunctionTok{\},}
    \DataTypeTok{"[json]"}\FunctionTok{:} \FunctionTok{\{}
        \DataTypeTok{"editor.defaultFormatter"}\FunctionTok{:} \StringTok{"vscode.json{-}language{-}features"}
    \FunctionTok{\},}
    
    \ErrorTok{//} \ErrorTok{==========================================}
    \ErrorTok{//} \ErrorTok{QUARTO}
    \ErrorTok{//} \ErrorTok{==========================================}
    \DataTypeTok{"quarto.visualEditor.spellingDictionary"}\FunctionTok{:} \StringTok{"es\_ES"}\FunctionTok{,}
    \DataTypeTok{"quarto.render.previewType"}\FunctionTok{:} \StringTok{"internal"}\FunctionTok{,}
    
    \ErrorTok{//} \ErrorTok{==========================================}
    \ErrorTok{//} \ErrorTok{POSITRON{-}SPECIFIC}
    \ErrorTok{//} \ErrorTok{==========================================}
    \DataTypeTok{"positron.assistant.enable"}\FunctionTok{:} \KeywordTok{true}\FunctionTok{,}
    \DataTypeTok{"positron.assistant.alwaysEnableApplyInEditorAction"}\FunctionTok{:} \KeywordTok{true}\FunctionTok{,}
    \DataTypeTok{"positron.assistant.alwaysIncludeCopilotTools"}\FunctionTok{:} \KeywordTok{true}\FunctionTok{,}
    \DataTypeTok{"positron.assistant.showTokenUsage.enable"}\FunctionTok{:} \KeywordTok{true}
\FunctionTok{\}}
\end{Highlighting}
\end{Shaded}

\section{Python en Positron}\label{python-en-positron}

\subsection{Prerequisitos}\label{prerequisitos}

\textbf{Versiones soportadas:} Python 3.9 - 3.13

\textbf{Gestores recomendados:}

\begin{itemize}
\tightlist
\item
  \texttt{uv} (recomendado)
\item
  \texttt{pyenv}
\item
  \texttt{conda}/\texttt{mamba}
\item
  \texttt{venv}
\end{itemize}

\subsection{Instalar Python}\label{instalar-python}

\textbf{Con uv (recomendado):}

\begin{Shaded}
\begin{Highlighting}[]
\CommentTok{\# Instalar uv}
\ExtensionTok{curl} \AttributeTok{{-}LsSf}\NormalTok{ https://astral.sh/uv/install.sh }\KeywordTok{|} \FunctionTok{sh}

\CommentTok{\# Instalar Python}
\ExtensionTok{uv}\NormalTok{ python install 3.12}

\CommentTok{\# Ver versiones instaladas}
\ExtensionTok{uv}\NormalTok{ python list}
\end{Highlighting}
\end{Shaded}

\textbf{Con pyenv:}

\begin{Shaded}
\begin{Highlighting}[]
\CommentTok{\# Instalar pyenv}
\ExtensionTok{curl}\NormalTok{ https://pyenv.run }\KeywordTok{|} \FunctionTok{bash}

\CommentTok{\# Instalar Python}
\ExtensionTok{pyenv}\NormalTok{ install 3.12.0}
\ExtensionTok{pyenv}\NormalTok{ global 3.12.0}
\end{Highlighting}
\end{Shaded}

\textbf{Con conda:}

\begin{Shaded}
\begin{Highlighting}[]
\CommentTok{\# Crear entorno}
\ExtensionTok{conda}\NormalTok{ create }\AttributeTok{{-}n}\NormalTok{ myenv python=3.12}
\ExtensionTok{conda}\NormalTok{ activate myenv}
\end{Highlighting}
\end{Shaded}

\subsection{Configurar Python en
Positron}\label{configurar-python-en-positron}

\textbf{Settings para Anaconda:}

\begin{Shaded}
\begin{Highlighting}[]
\FunctionTok{\{}
    \DataTypeTok{"python.defaultInterpreterPath"}\FunctionTok{:} \StringTok{"/home/achalmaedison/anaconda3/bin/python"}\FunctionTok{,}
    \DataTypeTok{"python.condaPath"}\FunctionTok{:} \StringTok{"/home/achalmaedison/anaconda3/bin/conda"}\FunctionTok{,}
    \DataTypeTok{"python.terminal.activateEnvironment"}\FunctionTok{:} \KeywordTok{true}\FunctionTok{,}
    \DataTypeTok{"python.terminal.activateEnvInCurrentTerminal"}\FunctionTok{:} \KeywordTok{true}
\FunctionTok{\}}
\end{Highlighting}
\end{Shaded}

\textbf{IPython Kernel:}

\begin{Shaded}
\begin{Highlighting}[]
\FunctionTok{\{}
    \ErrorTok{//} \ErrorTok{Usar} \ErrorTok{IPython} \ErrorTok{kernel} \ErrorTok{empaquetado} \ErrorTok{(recomendado)}
    \DataTypeTok{"python.useBundledIpykernel"}\FunctionTok{:} \KeywordTok{true}\FunctionTok{,}
    
    \ErrorTok{//} \ErrorTok{O} \ErrorTok{instalar} \ErrorTok{manualmente}
    \ErrorTok{//} \ErrorTok{python.useBundledIpykernel}\DataTypeTok{": false}
\DataTypeTok{\}}
\end{Highlighting}
\end{Shaded}

\subsection{Seleccionar Intérprete}\label{seleccionar-intuxe9rprete}

\textbf{Command Palette:}

\begin{verbatim}
Python: Select Interpreter
\end{verbatim}

\textbf{Desde Console:} Click en el selector de intérprete en la esquina
superior derecha

\subsection{Características Python}\label{caracteruxedsticas-python}

\textbf{Code Completion:}

\begin{itemize}
\tightlist
\item
  IntelliSense (Jedi)
\item
  Type hints
\item
  Import suggestions
\end{itemize}

\textbf{Type Checking:}

\begin{itemize}
\tightlist
\item
  Pyright integrado
\item
  Configuración en \texttt{pyrightconfig.json}
\end{itemize}

\textbf{Formatting:}

\begin{itemize}
\tightlist
\item
  Ruff (recomendado)
\item
  Black
\item
  autopep8
\end{itemize}

\textbf{Linting:}

\begin{itemize}
\tightlist
\item
  Ruff
\item
  Pylint
\item
  Flake8
\end{itemize}

\subsection{Configuración Python
Avanzada}\label{configuraciuxf3n-python-avanzada}

\begin{Shaded}
\begin{Highlighting}[]
\FunctionTok{\{}
    \ErrorTok{//} \ErrorTok{Intérprete}
    \DataTypeTok{"python.defaultInterpreterPath"}\FunctionTok{:} \StringTok{"/path/to/python"}\FunctionTok{,}
    
    \ErrorTok{//} \ErrorTok{IPython}
    \DataTypeTok{"python.useBundledIpykernel"}\FunctionTok{:} \KeywordTok{true}\FunctionTok{,}
    \DataTypeTok{"python.enableAutoReload"}\FunctionTok{:} \KeywordTok{true}\FunctionTok{,}
    
    \ErrorTok{//} \ErrorTok{Formatting}
    \DataTypeTok{"python.formatting.provider"}\FunctionTok{:} \StringTok{"none"}\FunctionTok{,}
    \DataTypeTok{"[python]"}\FunctionTok{:} \FunctionTok{\{}
        \DataTypeTok{"editor.defaultFormatter"}\FunctionTok{:} \StringTok{"charliermarsh.ruff"}\FunctionTok{,}
        \DataTypeTok{"editor.formatOnSave"}\FunctionTok{:} \KeywordTok{true}\FunctionTok{,}
        \DataTypeTok{"editor.codeActionsOnSave"}\FunctionTok{:} \FunctionTok{\{}
            \DataTypeTok{"source.organizeImports"}\FunctionTok{:} \StringTok{"explicit"}\FunctionTok{,}
            \DataTypeTok{"source.fixAll"}\FunctionTok{:} \StringTok{"explicit"}
        \FunctionTok{\}}
    \FunctionTok{\},}
    
    \ErrorTok{//} \ErrorTok{Linting}
    \DataTypeTok{"python.linting.enabled"}\FunctionTok{:} \KeywordTok{true}\FunctionTok{,}
    \DataTypeTok{"python.linting.ruffEnabled"}\FunctionTok{:} \KeywordTok{true}\FunctionTok{,}
    
    \ErrorTok{//} \ErrorTok{Terminal}
    \DataTypeTok{"python.terminal.activateEnvironment"}\FunctionTok{:} \KeywordTok{true}\FunctionTok{,}
    \DataTypeTok{"python.terminal.activateEnvInCurrentTerminal"}\FunctionTok{:} \KeywordTok{true}
\FunctionTok{\}}
\end{Highlighting}
\end{Shaded}

\subsection{Paquetes Esenciales}\label{paquetes-esenciales}

\begin{Shaded}
\begin{Highlighting}[]
\CommentTok{\# Data Science}
\ExtensionTok{pip}\NormalTok{ install numpy pandas matplotlib seaborn scikit{-}learn}

\CommentTok{\# Jupyter}
\ExtensionTok{pip}\NormalTok{ install ipykernel jupyter}

\CommentTok{\# Development}
\ExtensionTok{pip}\NormalTok{ install ruff black isort mypy pytest}

\CommentTok{\# Para Positron Console}
\ExtensionTok{pip}\NormalTok{ install ipython}
\end{Highlighting}
\end{Shaded}

\section{R en Positron}\label{r-en-positron}

\subsection{Prerequisitos}\label{prerequisitos-1}

\textbf{Versión requerida:} R 4.2 o superior

\subsection{Instalar R}\label{instalar-r}

\textbf{Linux:}

\begin{Shaded}
\begin{Highlighting}[]
\CommentTok{\# Ubuntu/Debian}
\FunctionTok{sudo}\NormalTok{ apt install r{-}base r{-}base{-}dev}

\CommentTok{\# O desde CRAN}
\CommentTok{\# https://cloud.r{-}project.org/}
\end{Highlighting}
\end{Shaded}

\textbf{macOS:}

\begin{Shaded}
\begin{Highlighting}[]
\CommentTok{\# Con Homebrew}
\ExtensionTok{brew}\NormalTok{ install r}

\CommentTok{\# O descargar desde CRAN}
\CommentTok{\# https://cloud.r{-}project.org/}
\end{Highlighting}
\end{Shaded}

\textbf{Con rig (múltiples versiones):}

\begin{Shaded}
\begin{Highlighting}[]
\CommentTok{\# Instalar rig}
\ExtensionTok{curl} \AttributeTok{{-}Ls}\NormalTok{ https://github.com/r{-}lib/rig/releases/latest/download/rig{-}linux{-}latest.tar.gz }\KeywordTok{|} \FunctionTok{sudo}\NormalTok{ tar xz }\AttributeTok{{-}C}\NormalTok{ /usr/local}

\CommentTok{\# Instalar R}
\ExtensionTok{rig}\NormalTok{ add release}
\ExtensionTok{rig}\NormalTok{ add devel}

\CommentTok{\# Cambiar versión}
\ExtensionTok{rig}\NormalTok{ default 4.3}
\end{Highlighting}
\end{Shaded}

\subsection{Paquetes R Esenciales para
Positron}\label{paquetes-r-esenciales-para-positron}

\begin{Shaded}
\begin{Highlighting}[]
\CommentTok{\# Paquetes core}
\FunctionTok{install.packages}\NormalTok{(}\FunctionTok{c}\NormalTok{(}
    \StringTok{"usethis"}\NormalTok{, }
    \StringTok{"cli"}\NormalTok{, }
    \StringTok{"crayon"}\NormalTok{, }
    \StringTok{"rlang"}\NormalTok{, }
    \StringTok{"roxygen2"}\NormalTok{, }
    \StringTok{"pkgload"}
\NormalTok{))}

\CommentTok{\# Con pak (recomendado)}
\NormalTok{pak}\SpecialCharTok{::}\FunctionTok{pak}\NormalTok{(}\FunctionTok{c}\NormalTok{(}
    \StringTok{"usethis"}\NormalTok{, }
    \StringTok{"cli"}\NormalTok{, }
    \StringTok{"crayon"}\NormalTok{, }
    \StringTok{"rlang"}\NormalTok{, }
    \StringTok{"roxygen2"}\NormalTok{, }
    \StringTok{"pkgload"}
\NormalTok{))}

\CommentTok{\# Tidyverse}
\FunctionTok{install.packages}\NormalTok{(}\StringTok{"tidyverse"}\NormalTok{)}

\CommentTok{\# Desarrollo}
\FunctionTok{install.packages}\NormalTok{(}\FunctionTok{c}\NormalTok{(}
    \StringTok{"devtools"}\NormalTok{,}
    \StringTok{"testthat"}\NormalTok{,}
    \StringTok{"lintr"}\NormalTok{,}
    \StringTok{"styler"}
\NormalTok{))}
\end{Highlighting}
\end{Shaded}

\subsection{Configurar R en Positron}\label{configurar-r-en-positron}

\textbf{Seleccionar intérprete R:}

\begin{verbatim}
Command Palette → R: Select Interpreter
\end{verbatim}

\textbf{Settings:}

\begin{Shaded}
\begin{Highlighting}[]
\FunctionTok{\{}
    \DataTypeTok{"r.alwaysUseActiveTerminal"}\FunctionTok{:} \KeywordTok{true}\FunctionTok{,}
    \DataTypeTok{"r.bracketedPaste"}\FunctionTok{:} \KeywordTok{true}\FunctionTok{,}
    \DataTypeTok{"r.plot.useHttpgd"}\FunctionTok{:} \KeywordTok{true}
\FunctionTok{\}}
\end{Highlighting}
\end{Shaded}

\subsection{IRKernel}\label{irkernel}

Positron maneja IRKernel automáticamente. Si necesitas configurarlo
manualmente:

\begin{Shaded}
\begin{Highlighting}[]
\CommentTok{\# Instalar IRKernel}
\FunctionTok{install.packages}\NormalTok{(}\StringTok{\textquotesingle{}IRkernel\textquotesingle{}}\NormalTok{)}

\CommentTok{\# Registrar con Jupyter}
\NormalTok{IRkernel}\SpecialCharTok{::}\FunctionTok{installspec}\NormalTok{(}\AttributeTok{user =} \ConstantTok{TRUE}\NormalTok{)}
\end{Highlighting}
\end{Shaded}

\subsection{Características R}\label{caracteruxedsticas-r}

\textbf{Code Completion:}

\begin{itemize}
\tightlist
\item
  Funciones
\item
  Paquetes
\item
  Variables
\item
  Archivos
\end{itemize}

\textbf{Help:}

\begin{itemize}
\tightlist
\item
  Documentación integrada
\item
  Ejemplos
\item
  Vignettes
\end{itemize}

\textbf{Debugging:}

\begin{itemize}
\tightlist
\item
  Breakpoints
\item
  Step through
\item
  Inspect variables
\end{itemize}

\subsection{Rtools (Windows)}\label{rtools-windows}

Si usas Windows y desarrollas paquetes:

\begin{Shaded}
\begin{Highlighting}[]
\CommentTok{\# Con rig}
\ExtensionTok{rig}\NormalTok{ system rtools add}

\CommentTok{\# O manual desde CRAN}
\CommentTok{\# https://cran.r{-}project.org/bin/windows/Rtools/}
\end{Highlighting}
\end{Shaded}

\section{Extensiones}\label{extensiones}

\subsection{Open VSX vs VS
Marketplace}\label{open-vsx-vs-vs-marketplace}

Positron usa \textbf{Open VSX} por razones de licencia. La mayoría de
extensiones populares están disponibles, pero algunas pueden faltar.

\subsection{Acceder a Extensiones}\label{acceder-a-extensiones}

\begin{verbatim}
View → Extensions (Ctrl+Shift+X)
\end{verbatim}

\subsection{Extensiones
Pre-instaladas}\label{extensiones-pre-instaladas}

\textbf{Built-in (core):}

\begin{itemize}
\tightlist
\item
  Python backend (posit.positron-python)
\item
  R backend (posit.air-vscode)
\end{itemize}

\textbf{Bootstrapped (auto-instaladas):}

\begin{itemize}
\tightlist
\item
  \texttt{charliermarsh.ruff}: Linter Python
\item
  \texttt{meta.pyrefly}: Python language server
\item
  \texttt{posit.shiny}: Soporte Shiny
\item
  \texttt{posit.publisher}: Publishing tool
\item
  \texttt{quarto.quarto}: Soporte Quarto
\end{itemize}

\subsection{Extensiones Recomendadas}\label{extensiones-recomendadas}

\textbf{Temas:}

\begin{verbatim}
- material-icon-theme (iconos)
- One Dark Pro
- GitHub Theme
- Dracula
- Nord
\end{verbatim}

\textbf{Productividad:}

\begin{verbatim}
- GitLens (si disponible)
- Better Comments
- Bookmarks
- Todo Tree
\end{verbatim}

\textbf{Formateo:}

\begin{verbatim}
- Prettier (esbenp.prettier-vscode)
- EditorConfig
\end{verbatim}

\textbf{Markdown/Docs:}

\begin{verbatim}
- Markdown All in One
- Markdown Preview Enhanced
\end{verbatim}

\textbf{LaTeX:}

\begin{verbatim}
- LaTeX Workshop (James-Yu.latex-workshop)
\end{verbatim}

\subsection{Instalar Extensión}\label{instalar-extensiuxf3n}

\textbf{Desde GUI:}

\begin{enumerate}
\def\labelenumi{\arabic{enumi}.}
\tightlist
\item
  Abrir Extensions (Ctrl+Shift+X)
\item
  Buscar extensión
\item
  Click en Install
\end{enumerate}

\textbf{Desde VSIX:}

\begin{verbatim}
Command Palette → Extensions: Install from VSIX
\end{verbatim}

\subsection{Gestionar Extensiones
Bootstrapped}\label{gestionar-extensiones-bootstrapped}

Si quieres deshabilitar una extensión bootstrapped:

\begin{Shaded}
\begin{Highlighting}[]
\FunctionTok{\{}
    \DataTypeTok{"extensions.allowed"}\FunctionTok{:} \FunctionTok{\{}
        \DataTypeTok{"posit.air{-}vscode"}\FunctionTok{:} \KeywordTok{false}\FunctionTok{,}  \ErrorTok{//} \ErrorTok{Deshabilitar} \ErrorTok{R} \ErrorTok{si} \ErrorTok{solo} \ErrorTok{usas} \ErrorTok{Python}
        \DataTypeTok{"*"}\FunctionTok{:} \KeywordTok{true}
    \FunctionTok{\}}
\FunctionTok{\}}
\end{Highlighting}
\end{Shaded}

\subsection{Extensiones No
Compatibles}\label{extensiones-no-compatibles}

\textbf{NO usar:}

\begin{itemize}
\tightlist
\item
  \texttt{ms-python.python} (Positron tiene soporte Python nativo)
\item
  \texttt{REditorSupport.r} (conflicto con posit.air-vscode)
\end{itemize}

\section{Temas y Personalización
Visual}\label{temas-y-personalizaciuxf3n-visual}

\subsection{Temas de Color}\label{temas-de-color}

\textbf{Cambiar tema:}

\begin{verbatim}
Command Palette → Preferences: Color Theme (Ctrl+K Ctrl+T)
\end{verbatim}

\textbf{Temas incluidos:}

\begin{itemize}
\tightlist
\item
  Default Positron Dark
\item
  Default Positron Light
\item
  Dark+ (heredado de VS Code)
\item
  Light+ (heredado de VS Code)
\end{itemize}

\textbf{Instalar más temas:}

\begin{verbatim}
Extensions → Buscar "theme"
\end{verbatim}

\textbf{Temas populares:}

\begin{itemize}
\tightlist
\item
  One Dark Pro
\item
  Material Theme
\item
  GitHub Theme
\item
  Dracula
\item
  Nord
\item
  Gruvbox
\end{itemize}

\textbf{En settings.json:}

\begin{Shaded}
\begin{Highlighting}[]
\FunctionTok{\{}
    \DataTypeTok{"workbench.colorTheme"}\FunctionTok{:} \StringTok{"One Dark Pro Mix"}
\FunctionTok{\}}
\end{Highlighting}
\end{Shaded}

\subsection{Temas de Iconos}\label{temas-de-iconos}

\begin{verbatim}
Command Palette → Preferences: File Icon Theme
\end{verbatim}

\textbf{Recomendados:}

\begin{itemize}
\tightlist
\item
  Material Icon Theme (recomendado)
\item
  Seti
\item
  Minimal
\end{itemize}

\begin{Shaded}
\begin{Highlighting}[]
\FunctionTok{\{}
    \DataTypeTok{"workbench.iconTheme"}\FunctionTok{:} \StringTok{"material{-}icon{-}theme"}
\FunctionTok{\}}
\end{Highlighting}
\end{Shaded}

\subsection{Customizar Colores}\label{customizar-colores}

\begin{Shaded}
\begin{Highlighting}[]
\FunctionTok{\{}
    \DataTypeTok{"workbench.colorCustomizations"}\FunctionTok{:} \FunctionTok{\{}
        \DataTypeTok{"[Default Positron Dark]"}\FunctionTok{:} \FunctionTok{\{}
            \DataTypeTok{"editor.background"}\FunctionTok{:} \StringTok{"\#1a1a1a"}\FunctionTok{,}
            \DataTypeTok{"sideBar.background"}\FunctionTok{:} \StringTok{"\#1e1e1e"}\FunctionTok{,}
            \DataTypeTok{"activityBar.background"}\FunctionTok{:} \StringTok{"\#252525"}
        \FunctionTok{\}}
    \FunctionTok{\},}
    \DataTypeTok{"editor.tokenColorCustomizations"}\FunctionTok{:} \FunctionTok{\{}
        \DataTypeTok{"[Default Positron Dark]"}\FunctionTok{:} \FunctionTok{\{}
            \DataTypeTok{"comments"}\FunctionTok{:} \StringTok{"\#6A9955"}\FunctionTok{,}
            \DataTypeTok{"strings"}\FunctionTok{:} \StringTok{"\#CE9178"}
        \FunctionTok{\}}
    \FunctionTok{\}}
\FunctionTok{\}}
\end{Highlighting}
\end{Shaded}

\subsection{Fuentes}\label{fuentes}

\textbf{Fuentes recomendadas con ligaduras:}

\begin{itemize}
\tightlist
\item
  JetBrains Mono
\item
  Fira Code
\item
  Cascadia Code
\item
  Hack Nerd Font
\item
  Victor Mono
\end{itemize}

\textbf{Configurar:}

\begin{Shaded}
\begin{Highlighting}[]
\FunctionTok{\{}
    \DataTypeTok{"editor.fontFamily"}\FunctionTok{:} \StringTok{"\textquotesingle{}JetBrains Mono\textquotesingle{}, \textquotesingle{}Fira Code\textquotesingle{}, monospace"}\FunctionTok{,}
    \DataTypeTok{"editor.fontLigatures"}\FunctionTok{:} \KeywordTok{true}\FunctionTok{,}
    \DataTypeTok{"editor.fontSize"}\FunctionTok{:} \DecValTok{13}\FunctionTok{,}
    \DataTypeTok{"terminal.integrated.fontFamily"}\FunctionTok{:} \StringTok{"\textquotesingle{}Hack Nerd Font\textquotesingle{}, monospace"}\FunctionTok{,}
    \DataTypeTok{"terminal.integrated.fontSize"}\FunctionTok{:} \DecValTok{12}
\FunctionTok{\}}
\end{Highlighting}
\end{Shaded}

\section{Terminal}\label{terminal}

\subsection{Terminal vs Console}\label{terminal-vs-console}

\textbf{Terminal:}

\begin{itemize}
\tightlist
\item
  Shell del sistema (bash, zsh, fish, etc.)
\item
  Comandos del sistema
\item
  Instalación de paquetes
\item
  Operaciones de archivos
\end{itemize}

\textbf{Console:}

\begin{itemize}
\tightlist
\item
  Intérprete interactivo (IPython/IRKernel)
\item
  Ejecución de código Python/R
\item
  Inspección de variables
\item
  Integración con Variables/Plots panes
\end{itemize}

\subsection{Configurar Terminal}\label{configurar-terminal}

\begin{Shaded}
\begin{Highlighting}[]
\FunctionTok{\{}
    \ErrorTok{//} \ErrorTok{Perfil} \ErrorTok{por} \ErrorTok{defecto}
    \DataTypeTok{"terminal.integrated.defaultProfile.linux"}\FunctionTok{:} \StringTok{"bash"}\FunctionTok{,}
    
    \ErrorTok{//} \ErrorTok{Fuente}
    \DataTypeTok{"terminal.integrated.fontFamily"}\FunctionTok{:} \StringTok{"\textquotesingle{}Hack Nerd Font\textquotesingle{}, monospace"}\FunctionTok{,}
    \DataTypeTok{"terminal.integrated.fontSize"}\FunctionTok{:} \DecValTok{12}\FunctionTok{,}
    
    \ErrorTok{//} \ErrorTok{Comportamiento}
    \DataTypeTok{"terminal.integrated.copyOnSelection"}\FunctionTok{:} \KeywordTok{true}\FunctionTok{,}
    \DataTypeTok{"terminal.integrated.cursorBlinking"}\FunctionTok{:} \KeywordTok{true}\FunctionTok{,}
    \DataTypeTok{"terminal.integrated.cursorStyle"}\FunctionTok{:} \StringTok{"line"}\FunctionTok{,}
    \DataTypeTok{"terminal.integrated.scrollback"}\FunctionTok{:} \DecValTok{10000}\FunctionTok{,}
    
    \ErrorTok{//} \ErrorTok{Entorno}
    \DataTypeTok{"terminal.integrated.inheritEnv"}\FunctionTok{:} \KeywordTok{false}
\FunctionTok{\}}
\end{Highlighting}
\end{Shaded}

\subsection{Perfiles de Terminal}\label{perfiles-de-terminal}

\begin{Shaded}
\begin{Highlighting}[]
\FunctionTok{\{}
    \DataTypeTok{"terminal.integrated.profiles.linux"}\FunctionTok{:} \FunctionTok{\{}
        \DataTypeTok{"bash"}\FunctionTok{:} \FunctionTok{\{}
            \DataTypeTok{"path"}\FunctionTok{:} \StringTok{"/bin/bash"}\FunctionTok{,}
            \DataTypeTok{"icon"}\FunctionTok{:} \StringTok{"terminal{-}bash"}
        \FunctionTok{\},}
        \DataTypeTok{"zsh"}\FunctionTok{:} \FunctionTok{\{}
            \DataTypeTok{"path"}\FunctionTok{:} \StringTok{"/usr/bin/zsh"}\FunctionTok{,}
            \DataTypeTok{"icon"}\FunctionTok{:} \StringTok{"terminal"}
        \FunctionTok{\},}
        \DataTypeTok{"fish"}\FunctionTok{:} \FunctionTok{\{}
            \DataTypeTok{"path"}\FunctionTok{:} \StringTok{"/usr/bin/fish"}\FunctionTok{,}
            \DataTypeTok{"icon"}\FunctionTok{:} \StringTok{"terminal"}
        \FunctionTok{\}}
    \FunctionTok{\}}
\FunctionTok{\}}
\end{Highlighting}
\end{Shaded}

\subsection{Atajos de Terminal}\label{atajos-de-terminal}

\begin{longtable}[]{@{}ll@{}}
\toprule\noalign{}
Atajo & Acción \\
\midrule\noalign{}
\endhead
\bottomrule\noalign{}
\endlastfoot
\texttt{Ctrl+\textasciigrave{}} & Toggle terminal \\
\texttt{Ctrl+Shift+\textasciigrave{}} & New terminal \\
\texttt{Ctrl+Shift+C} & Copy \\
\texttt{Ctrl+Shift+V} & Paste \\
\texttt{Ctrl+Shift+5} & Split terminal \\
\end{longtable}

\section{Console}\label{console}

\subsection{¿Qué es el Console?}\label{quuxe9-es-el-console}

El Console es una shell interactiva para Python (IPython) o R (IRKernel)
integrada directamente en Positron.

\textbf{Características:}

\begin{itemize}
\tightlist
\item
  Ejecución interactiva de código
\item
  Code completion
\item
  Syntax highlighting
\item
  Integración con Variables pane
\item
  Integración con Plots pane
\item
  Historial persistente
\end{itemize}

\subsection{Ejecutar Código en
Console}\label{ejecutar-cuxf3digo-en-console}

\textbf{Desde Editor:}

\begin{itemize}
\tightlist
\item
  \texttt{Ctrl+Enter}: Ejecutar línea/selección
\item
  \texttt{Ctrl+Alt+Home}: Ejecutar desde inicio hasta cursor
\item
  \texttt{Ctrl+Alt+End}: Ejecutar desde cursor hasta fin
\item
  \texttt{Ctrl+Shift+P}: Ejecutar archivo completo
\end{itemize}

\textbf{Directamente en Console:}

\begin{itemize}
\tightlist
\item
  Escribir código y presionar Enter
\end{itemize}

\subsection{Console Settings}\label{console-settings}

\begin{Shaded}
\begin{Highlighting}[]
\FunctionTok{\{}
    \ErrorTok{//} \ErrorTok{Python}
    \DataTypeTok{"python.useBundledIpykernel"}\FunctionTok{:} \KeywordTok{true}\FunctionTok{,}
    \DataTypeTok{"python.enableAutoReload"}\FunctionTok{:} \KeywordTok{true}\FunctionTok{,}
    
    \ErrorTok{//} \ErrorTok{R}
    \DataTypeTok{"r.alwaysUseActiveTerminal"}\FunctionTok{:} \KeywordTok{true}
\FunctionTok{\}}
\end{Highlighting}
\end{Shaded}

\subsection{Atajos Console}\label{atajos-console}

\begin{longtable}[]{@{}ll@{}}
\toprule\noalign{}
Atajo & Acción \\
\midrule\noalign{}
\endhead
\bottomrule\noalign{}
\endlastfoot
\texttt{Ctrl+K,\ F} & Enfocar Console \\
\texttt{Ctrl+L} & Limpiar Console \\
\texttt{Ctrl+Shift+0} & Reiniciar intérprete \\
\texttt{Ctrl+Up/Down} & Navegar historial \\
\end{longtable}

\section{Variables Pane}\label{variables-pane}

\subsection{¿Qué es Variables Pane?}\label{quuxe9-es-variables-pane}

Panel que muestra todas las variables en el entorno actual.

\textbf{Información mostrada:}

\begin{itemize}
\tightlist
\item
  Nombre
\item
  Tipo
\item
  Valor (para primitivos)
\item
  Tamaño (para colecciones)
\item
  Dimensiones (para dataframes)
\end{itemize}

\subsection{Usar Variables Pane}\label{usar-variables-pane}

\textbf{Abrir:}

\begin{verbatim}
View → Variables (Ctrl+K, V)
\end{verbatim}

\textbf{Operaciones:}

\begin{itemize}
\tightlist
\item
  Click en variable para expandir
\item
  Double-click para abrir en Data Explorer (si es dataframe)
\item
  Right-click para opciones (copy, delete, etc.)
\end{itemize}

\subsection{Filtrar Variables}\label{filtrar-variables}

Usa la barra de búsqueda en el panel para filtrar variables por nombre.

\section{Data Explorer}\label{data-explorer}

\subsection{¿Qué es Data Explorer?}\label{quuxe9-es-data-explorer}

Visualizador interactivo de dataframes (pandas, tibble, data.frame).

\textbf{Características:}

\begin{itemize}
\tightlist
\item
  Vista tabular de datos
\item
  Ordenamiento por columnas
\item
  Filtrado
\item
  Búsqueda
\item
  Estadísticas de columnas
\item
  Exportar a CSV/Excel
\end{itemize}

\subsection{Abrir Data Explorer}\label{abrir-data-explorer}

\textbf{Desde Variables Pane:} Double-click en un dataframe

\textbf{Desde Console:}

\begin{Shaded}
\begin{Highlighting}[]
\CommentTok{\# Python}
\ImportTok{import}\NormalTok{ pandas }\ImportTok{as}\NormalTok{ pd}
\NormalTok{df }\OperatorTok{=}\NormalTok{ pd.read\_csv(}\StringTok{\textquotesingle{}data.csv\textquotesingle{}}\NormalTok{)}
\CommentTok{\# Double{-}click en \textquotesingle{}df\textquotesingle{} en Variables pane}
\end{Highlighting}
\end{Shaded}

\begin{Shaded}
\begin{Highlighting}[]
\CommentTok{\# R}
\NormalTok{df }\OtherTok{\textless{}{-}} \FunctionTok{read.csv}\NormalTok{(}\StringTok{\textquotesingle{}data.csv\textquotesingle{}}\NormalTok{)}
\FunctionTok{View}\NormalTok{(df)}
\end{Highlighting}
\end{Shaded}

\textbf{Command Palette:}

\begin{verbatim}
Data Explorer: Open Data Explorer
\end{verbatim}

\subsection{Usar Data Explorer}\label{usar-data-explorer}

\textbf{Ordenar:} Click en header de columna

\textbf{Filtrar:} Click en icono de filtro en header

\textbf{Buscar:} Ctrl+F en Data Explorer

\textbf{Exportar:} Right-click → Export

\section{Plots Pane}\label{plots-pane}

\subsection{¿Qué es Plots Pane?}\label{quuxe9-es-plots-pane}

Panel para visualización de gráficos generados en Python o R.

\textbf{Características:}

\begin{itemize}
\tightlist
\item
  Historial de plots
\item
  Zoom
\item
  Pan
\item
  Exportar (PNG, SVG, PDF)
\item
  Tamaño ajustable
\end{itemize}

\subsection{Generar Plots}\label{generar-plots}

\textbf{Python (matplotlib):}

\begin{Shaded}
\begin{Highlighting}[]
\ImportTok{import}\NormalTok{ matplotlib.pyplot }\ImportTok{as}\NormalTok{ plt}
\ImportTok{import}\NormalTok{ numpy }\ImportTok{as}\NormalTok{ np}

\NormalTok{x }\OperatorTok{=}\NormalTok{ np.linspace(}\DecValTok{0}\NormalTok{, }\DecValTok{10}\NormalTok{, }\DecValTok{100}\NormalTok{)}
\NormalTok{y }\OperatorTok{=}\NormalTok{ np.sin(x)}

\NormalTok{plt.plot(x, y)}
\NormalTok{plt.title(}\StringTok{\textquotesingle{}Sin Wave\textquotesingle{}}\NormalTok{)}
\NormalTok{plt.show()}
\end{Highlighting}
\end{Shaded}

\textbf{Python (seaborn):}

\begin{Shaded}
\begin{Highlighting}[]
\ImportTok{import}\NormalTok{ seaborn }\ImportTok{as}\NormalTok{ sns}
\ImportTok{import}\NormalTok{ pandas }\ImportTok{as}\NormalTok{ pd}

\NormalTok{df }\OperatorTok{=}\NormalTok{ pd.read\_csv(}\StringTok{\textquotesingle{}data.csv\textquotesingle{}}\NormalTok{)}
\NormalTok{sns.scatterplot(data}\OperatorTok{=}\NormalTok{df, x}\OperatorTok{=}\StringTok{\textquotesingle{}x\textquotesingle{}}\NormalTok{, y}\OperatorTok{=}\StringTok{\textquotesingle{}y\textquotesingle{}}\NormalTok{, hue}\OperatorTok{=}\StringTok{\textquotesingle{}category\textquotesingle{}}\NormalTok{)}
\NormalTok{plt.show()}
\end{Highlighting}
\end{Shaded}

\textbf{R (ggplot2):}

\begin{Shaded}
\begin{Highlighting}[]
\FunctionTok{library}\NormalTok{(ggplot2)}

\FunctionTok{ggplot}\NormalTok{(mtcars, }\FunctionTok{aes}\NormalTok{(}\AttributeTok{x =}\NormalTok{ mpg, }\AttributeTok{y =}\NormalTok{ hp)) }\SpecialCharTok{+}
  \FunctionTok{geom\_point}\NormalTok{() }\SpecialCharTok{+}
  \FunctionTok{theme\_minimal}\NormalTok{()}
\end{Highlighting}
\end{Shaded}

\subsection{Exportar Plots}\label{exportar-plots}

\textbf{Desde Plots Pane:}

\begin{enumerate}
\def\labelenumi{\arabic{enumi}.}
\tightlist
\item
  Hover sobre el plot
\item
  Click en icono de exportar
\item
  Elegir formato (PNG, SVG, PDF)
\item
  Especificar ubicación
\end{enumerate}

\textbf{Desde código:}

\begin{Shaded}
\begin{Highlighting}[]
\CommentTok{\# Python}
\NormalTok{plt.savefig(}\StringTok{\textquotesingle{}plot.png\textquotesingle{}}\NormalTok{, dpi}\OperatorTok{=}\DecValTok{300}\NormalTok{)}
\end{Highlighting}
\end{Shaded}

\begin{Shaded}
\begin{Highlighting}[]
\CommentTok{\# R}
\FunctionTok{ggsave}\NormalTok{(}\StringTok{\textquotesingle{}plot.png\textquotesingle{}}\NormalTok{, }\AttributeTok{width =} \DecValTok{8}\NormalTok{, }\AttributeTok{height =} \DecValTok{6}\NormalTok{, }\AttributeTok{dpi =} \DecValTok{300}\NormalTok{)}
\end{Highlighting}
\end{Shaded}

\section{Help Pane}\label{help-pane}

\subsection{¿Qué es Help Pane?}\label{quuxe9-es-help-pane}

Panel de documentación contextual integrado.

\textbf{Características:}

\begin{itemize}
\tightlist
\item
  Documentación de funciones
\item
  Argumentos
\item
  Ejemplos
\item
  Links a vignettes/tutorials
\item
  Búsqueda
\end{itemize}

\subsection{Obtener Ayuda}\label{obtener-ayuda}

\textbf{Atajo F1:} Coloca cursor sobre función y presiona F1

\textbf{Desde Console:}

\begin{Shaded}
\begin{Highlighting}[]
\CommentTok{\# Python}
\BuiltInTok{help}\NormalTok{(}\BuiltInTok{print}\NormalTok{)}
\NormalTok{?}\BuiltInTok{print}
\end{Highlighting}
\end{Shaded}

\begin{Shaded}
\begin{Highlighting}[]
\CommentTok{\# R}
\NormalTok{?mean}
\FunctionTok{help}\NormalTok{(mean)}
\end{Highlighting}
\end{Shaded}

\textbf{Command Palette:}

\begin{verbatim}
Help: Show Help for Symbol Under Cursor
\end{verbatim}

\section{Quarto}\label{quarto}

\subsection{¿Qué es Quarto?}\label{quuxe9-es-quarto}

Sistema de publicación científica y técnica de código abierto.

\textbf{Formatos de salida:}

\begin{itemize}
\tightlist
\item
  HTML
\item
  PDF
\item
  Word
\item
  Reveal.js (presentations)
\item
  Dashboards
\item
  Websites
\item
  Books
\end{itemize}

\subsection{Configurar Quarto}\label{configurar-quarto}

\textbf{Instalar Quarto:}

\begin{Shaded}
\begin{Highlighting}[]
\CommentTok{\# Linux}
\FunctionTok{sudo}\NormalTok{ curl }\AttributeTok{{-}LO}\NormalTok{ https://quarto.org/download/latest/quarto{-}linux{-}amd64.deb}
\FunctionTok{sudo}\NormalTok{ dpkg }\AttributeTok{{-}i}\NormalTok{ quarto{-}linux{-}amd64.deb}

\CommentTok{\# macOS}
\ExtensionTok{brew}\NormalTok{ install quarto}

\CommentTok{\# Verificar}
\ExtensionTok{quarto} \AttributeTok{{-}{-}version}
\end{Highlighting}
\end{Shaded}

\textbf{Settings:}

\begin{Shaded}
\begin{Highlighting}[]
\FunctionTok{\{}
    \DataTypeTok{"quarto.visualEditor.spellingDictionary"}\FunctionTok{:} \StringTok{"es\_ES"}\FunctionTok{,}
    \DataTypeTok{"quarto.render.previewType"}\FunctionTok{:} \StringTok{"internal"}\FunctionTok{,}
    \DataTypeTok{"[quarto]"}\FunctionTok{:} \FunctionTok{\{}
        \DataTypeTok{"editor.defaultFormatter"}\FunctionTok{:} \StringTok{"quarto.quarto"}\FunctionTok{,}
        \DataTypeTok{"editor.wordWrap"}\FunctionTok{:} \StringTok{"on"}
    \FunctionTok{\}}
\FunctionTok{\}}
\end{Highlighting}
\end{Shaded}

\subsection{Crear Documento Quarto}\label{crear-documento-quarto}

\textbf{Nuevo archivo:}

\begin{verbatim}
File → New File → Quarto Document
\end{verbatim}

\textbf{Estructura básica (.qmd):}

\begin{Shaded}
\begin{Highlighting}[]
\PreprocessorTok{{-}{-}{-}}
\FunctionTok{title}\KeywordTok{:}\AttributeTok{ }\StringTok{"Mi Documento"}
\FunctionTok{author}\KeywordTok{:}\AttributeTok{ }\StringTok{"Tu Nombre"}
\FunctionTok{format}\KeywordTok{:}\AttributeTok{ html}
\FunctionTok{execute}\KeywordTok{:}
\AttributeTok{  }\FunctionTok{echo}\KeywordTok{:}\AttributeTok{ }\CharTok{true}
\PreprocessorTok{{-}{-}{-}}

\CommentTok{\# Introducción}

\AttributeTok{Este es un documento Quarto.}

\AttributeTok{:}\FunctionTok{:}\KeywordTok{:}\AttributeTok{ }\KeywordTok{\{}\AttributeTok{.cell}\KeywordTok{\}}
\AttributeTok{\textasciigrave{}\textasciigrave{}\textasciigrave{} \{.python .cell{-}code\}}
\AttributeTok{import pandas as pd}
\AttributeTok{import matplotlib.pyplot as plt}

\AttributeTok{df = pd.read\_csv(\textquotesingle{}data.csv\textquotesingle{})}
\AttributeTok{df.head()}
\end{Highlighting}
\end{Shaded}

:::

\section{Análisis}\label{anuxe1lisis}

\begin{Shaded}
\begin{Highlighting}[]
\NormalTok{plt.plot(df[}\StringTok{\textquotesingle{}x\textquotesingle{}}\NormalTok{], df[}\StringTok{\textquotesingle{}y\textquotesingle{}}\NormalTok{])}
\NormalTok{plt.show()}
\end{Highlighting}
\end{Shaded}

\subsection{Renderizar}\label{renderizar}

\textbf{Command Palette:}

\begin{verbatim}
Quarto: Render Document
\end{verbatim}

\textbf{Atajo:}

\begin{verbatim}
Ctrl+Shift+K
\end{verbatim}

\textbf{Terminal:}

\begin{Shaded}
\begin{Highlighting}[]
\ExtensionTok{quarto}\NormalTok{ render document.qmd}
\ExtensionTok{quarto}\NormalTok{ render document.qmd }\AttributeTok{{-}{-}to}\NormalTok{ pdf}
\end{Highlighting}
\end{Shaded}

\subsection{Preview}\label{preview}

\textbf{Live preview:}

\begin{verbatim}
Quarto: Preview
\end{verbatim}

\subsection{Características Quarto en
Positron}\label{caracteruxedsticas-quarto-en-positron}

\begin{itemize}
\tightlist
\item
  Syntax highlighting
\item
  Code execution inline
\item
  Preview live
\item
  Autocompletado YAML
\item
  Snippets
\end{itemize}

\section{Jupyter Notebooks}\label{jupyter-notebooks}

\subsection{Soporte Jupyter}\label{soporte-jupyter}

Positron tiene soporte nativo para Jupyter notebooks (.ipynb).

\textbf{Características:}

\begin{itemize}
\tightlist
\item
  Ejecución de celdas
\item
  Output inline
\item
  Plots inline
\item
  Variables en Variables pane
\item
  Code completion
\end{itemize}

\subsection{Crear Notebook}\label{crear-notebook}

\begin{verbatim}
File → New File → Jupyter Notebook
\end{verbatim}

\subsection{Seleccionar Kernel}\label{seleccionar-kernel}

Click en ``Select Kernel'' en la esquina superior derecha

\subsection{Ejecutar Celdas}\label{ejecutar-celdas}

\textbf{Atajos:}

\begin{itemize}
\tightlist
\item
  \texttt{Shift+Enter}: Ejecutar celda y avanzar
\item
  \texttt{Ctrl+Enter}: Ejecutar celda
\item
  \texttt{Alt+Enter}: Ejecutar celda e insertar abajo
\end{itemize}

\subsection{Convertir entre Quarto y
Jupyter}\label{convertir-entre-quarto-y-jupyter}

\textbf{Quarto → Jupyter:}

\begin{Shaded}
\begin{Highlighting}[]
\ExtensionTok{quarto}\NormalTok{ convert document.qmd}
\end{Highlighting}
\end{Shaded}

\textbf{Jupyter → Quarto:}

\begin{Shaded}
\begin{Highlighting}[]
\ExtensionTok{quarto}\NormalTok{ convert notebook.ipynb}
\end{Highlighting}
\end{Shaded}

\section{Git y Control de Versiones}\label{git-y-control-de-versiones}

\subsection{Configurar Git}\label{configurar-git}

\textbf{Primera vez:}

\begin{Shaded}
\begin{Highlighting}[]
\FunctionTok{git}\NormalTok{ config }\AttributeTok{{-}{-}global}\NormalTok{ user.name }\StringTok{"Tu Nombre"}
\FunctionTok{git}\NormalTok{ config }\AttributeTok{{-}{-}global}\NormalTok{ user.email }\StringTok{"tu@email.com"}
\end{Highlighting}
\end{Shaded}

\textbf{Settings Positron:}

\begin{Shaded}
\begin{Highlighting}[]
\FunctionTok{\{}
    \DataTypeTok{"git.enableSmartCommit"}\FunctionTok{:} \KeywordTok{true}\FunctionTok{,}
    \DataTypeTok{"git.autofetch"}\FunctionTok{:} \KeywordTok{true}\FunctionTok{,}
    \DataTypeTok{"git.confirmSync"}\FunctionTok{:} \KeywordTok{false}\FunctionTok{,}
    \DataTypeTok{"git.openRepositoryInParentFolders"}\FunctionTok{:} \StringTok{"always"}
\FunctionTok{\}}
\end{Highlighting}
\end{Shaded}

\subsection{Source Control Panel}\label{source-control-panel}

\textbf{Abrir:}

\begin{verbatim}
View → Source Control (Ctrl+Shift+G)
\end{verbatim}

\textbf{Operaciones:}

\begin{itemize}
\tightlist
\item
  Ver cambios
\item
  Stage/Unstage
\item
  Commit
\item
  Push/Pull
\item
  Ver historial
\item
  Crear branches
\item
  Merge
\end{itemize}

\subsection{Workflow Git}\label{workflow-git}

\begin{verbatim}
1. Hacer cambios en archivos
2. Ver cambios en Source Control panel
3. Stage cambios (click en +)
4. Escribir mensaje de commit
5. Click en ✓ (commit)
6. Sync changes (push/pull)
\end{verbatim}

\subsection{GitLens (si disponible)}\label{gitlens-si-disponible}

Instalar desde Extensions si está disponible en Open VSX.

\section{AI Tools}\label{ai-tools}

\subsection{Positron Assistant}\label{positron-assistant}

Asistente AI integrado para ayudar con código.

\textbf{Habilitar:}

\begin{Shaded}
\begin{Highlighting}[]
\FunctionTok{\{}
    \DataTypeTok{"positron.assistant.enable"}\FunctionTok{:} \KeywordTok{true}\FunctionTok{,}
    \DataTypeTok{"positron.assistant.alwaysEnableApplyInEditorAction"}\FunctionTok{:} \KeywordTok{true}\FunctionTok{,}
    \DataTypeTok{"positron.assistant.showTokenUsage.enable"}\FunctionTok{:} \KeywordTok{true}
\FunctionTok{\}}
\end{Highlighting}
\end{Shaded}

\textbf{Usar:}

\begin{verbatim}
Command Palette → Positron: Open Assistant
\end{verbatim}

\textbf{Características:}

\begin{itemize}
\tightlist
\item
  Generar código
\item
  Explicar código
\item
  Refactorizar
\item
  Documentar
\item
  Debuggear
\end{itemize}

\subsection{Databot}\label{databot}

Asistente especializado para análisis de datos.

\textbf{Usar:}

\begin{verbatim}
Command Palette → Positron: Open Databot
\end{verbatim}

\textbf{Ejemplos de prompts:}

\begin{verbatim}
"Crear un gráfico de dispersión de x vs y"
"Agrupar por categoría y calcular media"
"Filtrar datos donde columna > 100"
"Crear modelo lineal de y ~ x"
\end{verbatim}

\section{Configuración Minimalista}\label{configuraciuxf3n-minimalista}

\subsection{Objetivos}\label{objetivos}

\begin{itemize}
\tightlist
\item
  UI limpia
\item
  Sin distracciones
\item
  Foco en código
\item
  Performance óptima
\end{itemize}

\subsection{Settings Minimalistas}\label{settings-minimalistas}

\begin{Shaded}
\begin{Highlighting}[]
\FunctionTok{\{}
    \ErrorTok{//} \ErrorTok{==========================================}
    \ErrorTok{//} \ErrorTok{MINIMALISTA}
    \ErrorTok{//} \ErrorTok{==========================================}
    
    \ErrorTok{//} \ErrorTok{Editor} \ErrorTok{limpio}
    \DataTypeTok{"editor.minimap.enabled"}\FunctionTok{:} \KeywordTok{false}\FunctionTok{,}
    \DataTypeTok{"editor.scrollbar.vertical"}\FunctionTok{:} \StringTok{"hidden"}\FunctionTok{,}
    \DataTypeTok{"editor.scrollbar.horizontal"}\FunctionTok{:} \StringTok{"hidden"}\FunctionTok{,}
    \DataTypeTok{"editor.lineNumbers"}\FunctionTok{:} \StringTok{"off"}\FunctionTok{,}  \ErrorTok{//} \ErrorTok{o} \DataTypeTok{"relative"}
    \DataTypeTok{"editor.glyphMargin"}\FunctionTok{:} \KeywordTok{false}\FunctionTok{,}
    \DataTypeTok{"editor.folding"}\FunctionTok{:} \KeywordTok{false}\FunctionTok{,}
    \DataTypeTok{"editor.guides.indentation"}\FunctionTok{:} \KeywordTok{false}\FunctionTok{,}
    \DataTypeTok{"editor.renderWhitespace"}\FunctionTok{:} \StringTok{"none"}\FunctionTok{,}
    \DataTypeTok{"editor.renderLineHighlight"}\FunctionTok{:} \StringTok{"none"}\FunctionTok{,}
    \DataTypeTok{"editor.overviewRulerBorder"}\FunctionTok{:} \KeywordTok{false}\FunctionTok{,}
    
    \ErrorTok{//} \ErrorTok{Workbench} \ErrorTok{limpio}
    \DataTypeTok{"workbench.statusBar.visible"}\FunctionTok{:} \KeywordTok{false}\FunctionTok{,}
    \DataTypeTok{"workbench.activityBar.location"}\FunctionTok{:} \StringTok{"hidden"}\FunctionTok{,}  \ErrorTok{//} \ErrorTok{o} \DataTypeTok{"bottom"}
    \DataTypeTok{"workbench.editor.showTabs"}\FunctionTok{:} \StringTok{"single"}\FunctionTok{,}
    \DataTypeTok{"workbench.editor.tabSizing"}\FunctionTok{:} \StringTok{"shrink"}\FunctionTok{,}
    \DataTypeTok{"breadcrumbs.enabled"}\FunctionTok{:} \KeywordTok{false}\FunctionTok{,}
    \DataTypeTok{"workbench.sideBar.location"}\FunctionTok{:} \StringTok{"right"}\FunctionTok{,}
    
    \ErrorTok{//} \ErrorTok{Terminal} \ErrorTok{limpio}
    \DataTypeTok{"terminal.integrated.tabs.enabled"}\FunctionTok{:} \KeywordTok{false}\FunctionTok{,}
    
    \ErrorTok{//} \ErrorTok{Tema} \ErrorTok{minimalista}
    \DataTypeTok{"workbench.colorTheme"}\FunctionTok{:} \StringTok{"Default Positron Dark"}\FunctionTok{,}
    \DataTypeTok{"workbench.iconTheme"}\FunctionTok{:} \StringTok{"minimal"}\FunctionTok{,}
    
    \ErrorTok{//} \ErrorTok{Zen} \ErrorTok{Mode} \ErrorTok{(Ctrl+K} \ErrorTok{Z)}
    \DataTypeTok{"zenMode.centerLayout"}\FunctionTok{:} \KeywordTok{true}\FunctionTok{,}
    \DataTypeTok{"zenMode.fullScreen"}\FunctionTok{:} \KeywordTok{false}\FunctionTok{,}
    \DataTypeTok{"zenMode.hideLineNumbers"}\FunctionTok{:} \KeywordTok{false}\FunctionTok{,}
    \DataTypeTok{"zenMode.hideTabs"}\FunctionTok{:} \KeywordTok{true}\FunctionTok{,}
    \DataTypeTok{"zenMode.hideStatusBar"}\FunctionTok{:} \KeywordTok{true}\FunctionTok{,}
    
    \ErrorTok{//} \ErrorTok{Resto} \ErrorTok{de} \ErrorTok{settings} \ErrorTok{esenciales}
    \DataTypeTok{"editor.fontSize"}\FunctionTok{:} \DecValTok{14}\FunctionTok{,}
    \DataTypeTok{"editor.fontFamily"}\FunctionTok{:} \StringTok{"\textquotesingle{}JetBrains Mono\textquotesingle{}, monospace"}\FunctionTok{,}
    \DataTypeTok{"editor.fontLigatures"}\FunctionTok{:} \KeywordTok{true}\FunctionTok{,}
    \DataTypeTok{"editor.formatOnSave"}\FunctionTok{:} \KeywordTok{true}\FunctionTok{,}
    \DataTypeTok{"files.autoSave"}\FunctionTok{:} \StringTok{"onFocusChange"}
\FunctionTok{\}}
\end{Highlighting}
\end{Shaded}

\subsection{Zen Mode}\label{zen-mode}

\textbf{Activar:}

\begin{verbatim}
Command Palette → View: Toggle Zen Mode (Ctrl+K Z)
\end{verbatim}

\textbf{O personalizar:}

\begin{Shaded}
\begin{Highlighting}[]
\FunctionTok{\{}
    \DataTypeTok{"zenMode.centerLayout"}\FunctionTok{:} \KeywordTok{true}\FunctionTok{,}
    \DataTypeTok{"zenMode.fullScreen"}\FunctionTok{:} \KeywordTok{false}\FunctionTok{,}
    \DataTypeTok{"zenMode.hideLineNumbers"}\FunctionTok{:} \KeywordTok{false}\FunctionTok{,}
    \DataTypeTok{"zenMode.hideTabs"}\FunctionTok{:} \KeywordTok{true}\FunctionTok{,}
    \DataTypeTok{"zenMode.hideStatusBar"}\FunctionTok{:} \KeywordTok{true}\FunctionTok{,}
    \DataTypeTok{"zenMode.hideActivityBar"}\FunctionTok{:} \KeywordTok{true}\FunctionTok{,}
    \DataTypeTok{"zenMode.restore"}\FunctionTok{:} \KeywordTok{true}
\FunctionTok{\}}
\end{Highlighting}
\end{Shaded}

\section{Mi Configuración
Específica}\label{mi-configuraciuxf3n-especuxedfica}

Esta es mi configuración de VS Code adaptada para Positron:

\begin{Shaded}
\begin{Highlighting}[]
\FunctionTok{\{}
    \ErrorTok{//} \ErrorTok{==========================================}
    \ErrorTok{//} \ErrorTok{PYTHON} \ErrorTok{(ANACONDA)}
    \ErrorTok{//} \ErrorTok{==========================================}
    \DataTypeTok{"python.defaultInterpreterPath"}\FunctionTok{:} \StringTok{"/home/achalmaedison/anaconda3/bin/python"}\FunctionTok{,}
    \DataTypeTok{"python.condaPath"}\FunctionTok{:} \StringTok{"/home/achalmaedison/anaconda3/bin/conda"}\FunctionTok{,}
    \DataTypeTok{"python.terminal.activateEnvironment"}\FunctionTok{:} \KeywordTok{true}\FunctionTok{,}
    \DataTypeTok{"python.terminal.activateEnvInCurrentTerminal"}\FunctionTok{:} \KeywordTok{true}\FunctionTok{,}
    \DataTypeTok{"python.useBundledIpykernel"}\FunctionTok{:} \KeywordTok{true}\FunctionTok{,}
    \DataTypeTok{"python.enableAutoReload"}\FunctionTok{:} \KeywordTok{true}\FunctionTok{,}
    
    \ErrorTok{//} \ErrorTok{==========================================}
    \ErrorTok{//} \ErrorTok{TERMINAL}
    \ErrorTok{//} \ErrorTok{==========================================}
    \DataTypeTok{"terminal.integrated.inheritEnv"}\FunctionTok{:} \KeywordTok{false}\FunctionTok{,}
    \DataTypeTok{"terminal.integrated.fontSize"}\FunctionTok{:} \DecValTok{12}\FunctionTok{,}
    \DataTypeTok{"terminal.integrated.fontFamily"}\FunctionTok{:} \StringTok{"Hack Nerd Font"}\FunctionTok{,}
    \DataTypeTok{"terminal.integrated.copyOnSelection"}\FunctionTok{:} \KeywordTok{true}\FunctionTok{,}
    \DataTypeTok{"terminal.integrated.cursorBlinking"}\FunctionTok{:} \KeywordTok{true}\FunctionTok{,}
    \DataTypeTok{"terminal.integrated.cursorStyle"}\FunctionTok{:} \StringTok{"underline"}\FunctionTok{,}
    \DataTypeTok{"terminal.integrated.cursorStyleInactive"}\FunctionTok{:} \StringTok{"line"}\FunctionTok{,}
    \DataTypeTok{"terminal.integrated.defaultProfile.linux"}\FunctionTok{:} \StringTok{"fish"}\FunctionTok{,}  \ErrorTok{//} \ErrorTok{Cambiado} \ErrorTok{de} \ErrorTok{zsh} \ErrorTok{a} \ErrorTok{fish}
    \DataTypeTok{"terminal.integrated.enablePersistentSessions"}\FunctionTok{:} \KeywordTok{false}\FunctionTok{,}
    \DataTypeTok{"terminal.integrated.suggest.enabled"}\FunctionTok{:} \KeywordTok{true}\FunctionTok{,}
    \DataTypeTok{"terminal.integrated.accessibleViewFocusOnCommandExecution"}\FunctionTok{:} \KeywordTok{true}\FunctionTok{,}
    
    \ErrorTok{//} \ErrorTok{==========================================}
    \ErrorTok{//} \ErrorTok{EDITOR}
    \ErrorTok{//} \ErrorTok{==========================================}
    \DataTypeTok{"editor.fontSize"}\FunctionTok{:} \DecValTok{13}\FunctionTok{,}
    \DataTypeTok{"editor.fontFamily"}\FunctionTok{:} \StringTok{"\textquotesingle{}Hack Nerd Font\textquotesingle{}, \textquotesingle{}JetBrains Mono\textquotesingle{}, monospace"}\FunctionTok{,}
    \DataTypeTok{"editor.fontLigatures"}\FunctionTok{:} \KeywordTok{true}\FunctionTok{,}
    \DataTypeTok{"editor.minimap.enabled"}\FunctionTok{:} \KeywordTok{false}\FunctionTok{,}
    \DataTypeTok{"editor.minimap.autohide"}\FunctionTok{:} \KeywordTok{true}\FunctionTok{,}
    \DataTypeTok{"editor.scrollbar.horizontal"}\FunctionTok{:} \StringTok{"hidden"}\FunctionTok{,}
    \DataTypeTok{"editor.scrollbar.vertical"}\FunctionTok{:} \StringTok{"hidden"}\FunctionTok{,}
    \DataTypeTok{"editor.overviewRulerBorder"}\FunctionTok{:} \KeywordTok{false}\FunctionTok{,}
    \DataTypeTok{"editor.formatOnSave"}\FunctionTok{:} \KeywordTok{true}\FunctionTok{,}
    \DataTypeTok{"editor.formatOnPaste"}\FunctionTok{:} \KeywordTok{true}\FunctionTok{,}
    \DataTypeTok{"editor.formatOnType"}\FunctionTok{:} \KeywordTok{true}\FunctionTok{,}
    \DataTypeTok{"editor.glyphMargin"}\FunctionTok{:} \KeywordTok{false}\FunctionTok{,}
    \DataTypeTok{"editor.linkedEditing"}\FunctionTok{:} \KeywordTok{true}\FunctionTok{,}
    \DataTypeTok{"editor.cursorBlinking"}\FunctionTok{:} \StringTok{"expand"}\FunctionTok{,}
    \DataTypeTok{"editor.cursorSmoothCaretAnimation"}\FunctionTok{:} \StringTok{"on"}\FunctionTok{,}
    \DataTypeTok{"editor.guides.indentation"}\FunctionTok{:} \KeywordTok{false}\FunctionTok{,}
    \DataTypeTok{"editor.matchBrackets"}\FunctionTok{:} \StringTok{"never"}\FunctionTok{,}
    \DataTypeTok{"editor.suggest.preview"}\FunctionTok{:} \KeywordTok{true}\FunctionTok{,}
    \DataTypeTok{"editor.inlineSuggest.syntaxHighlightingEnabled"}\FunctionTok{:} \KeywordTok{true}\FunctionTok{,}
    
    \ErrorTok{//} \ErrorTok{==========================================}
    \ErrorTok{//} \ErrorTok{WORKBENCH}
    \ErrorTok{//} \ErrorTok{==========================================}
    \DataTypeTok{"workbench.colorTheme"}\FunctionTok{:} \StringTok{"One Dark Pro Mix"}\FunctionTok{,}  \ErrorTok{//} \ErrorTok{Tu} \ErrorTok{tema} \ErrorTok{favorito}
    \DataTypeTok{"workbench.iconTheme"}\FunctionTok{:} \StringTok{"material{-}icon{-}theme"}\FunctionTok{,}
    \DataTypeTok{"workbench.productIconTheme"}\FunctionTok{:} \StringTok{"fluent{-}icons"}\FunctionTok{,}
    \DataTypeTok{"workbench.activityBar.location"}\FunctionTok{:} \StringTok{"bottom"}\FunctionTok{,}  \ErrorTok{//} \ErrorTok{Mantener} \ErrorTok{en} \ErrorTok{bottom}
    \DataTypeTok{"workbench.sideBar.location"}\FunctionTok{:} \StringTok{"right"}\FunctionTok{,}
    \DataTypeTok{"workbench.statusBar.visible"}\FunctionTok{:} \KeywordTok{false}\FunctionTok{,}  \ErrorTok{//} \ErrorTok{Oculto} \ErrorTok{para} \ErrorTok{minimalista}
    \DataTypeTok{"workbench.editor.showTabs"}\FunctionTok{:} \StringTok{"single"}\FunctionTok{,}
    \DataTypeTok{"workbench.editor.tabSizing"}\FunctionTok{:} \StringTok{"fixed"}\FunctionTok{,}
    \DataTypeTok{"workbench.editor.highlightModifiedTabs"}\FunctionTok{:} \KeywordTok{true}\FunctionTok{,}
    \DataTypeTok{"workbench.editor.limit.enabled"}\FunctionTok{:} \KeywordTok{true}\FunctionTok{,}
    \DataTypeTok{"workbench.editor.empty.hint"}\FunctionTok{:} \StringTok{"hidden"}\FunctionTok{,}
    \DataTypeTok{"workbench.editor.wrapTabs"}\FunctionTok{:} \KeywordTok{true}\FunctionTok{,}
    \DataTypeTok{"workbench.editor.centeredLayoutFixedWidth"}\FunctionTok{:} \KeywordTok{true}\FunctionTok{,}
    \DataTypeTok{"window.customTitleBarVisibility"}\FunctionTok{:} \StringTok{"auto"}\FunctionTok{,}
    \DataTypeTok{"breadcrumbs.enabled"}\FunctionTok{:} \KeywordTok{false}\FunctionTok{,}
    
    \ErrorTok{//} \ErrorTok{==========================================}
    \ErrorTok{//} \ErrorTok{FILES}
    \ErrorTok{//} \ErrorTok{==========================================}
    \DataTypeTok{"files.autoSave"}\FunctionTok{:} \StringTok{"onFocusChange"}\FunctionTok{,}
    \DataTypeTok{"files.associations"}\FunctionTok{:} \FunctionTok{\{}
        \DataTypeTok{"renv.lock"}\FunctionTok{:} \StringTok{"json"}\FunctionTok{,}
        \DataTypeTok{"*.qmd"}\FunctionTok{:} \StringTok{"quarto"}\FunctionTok{,}
        \DataTypeTok{"*.Rmd"}\FunctionTok{:} \StringTok{"rmarkdown"}
    \FunctionTok{\},}
    
    \ErrorTok{//} \ErrorTok{==========================================}
    \ErrorTok{//} \ErrorTok{GIT}
    \ErrorTok{//} \ErrorTok{==========================================}
    \DataTypeTok{"git.enableSmartCommit"}\FunctionTok{:} \KeywordTok{true}\FunctionTok{,}
    \DataTypeTok{"git.autofetch"}\FunctionTok{:} \KeywordTok{true}\FunctionTok{,}
    \DataTypeTok{"git.confirmSync"}\FunctionTok{:} \KeywordTok{false}\FunctionTok{,}
    \DataTypeTok{"git.openRepositoryInParentFolders"}\FunctionTok{:} \StringTok{"always"}\FunctionTok{,}
    
    \ErrorTok{//} \ErrorTok{==========================================}
    \ErrorTok{//} \ErrorTok{DIFF} \ErrorTok{EDITOR}
    \ErrorTok{//} \ErrorTok{==========================================}
    \DataTypeTok{"diffEditor.ignoreTrimWhitespace"}\FunctionTok{:} \KeywordTok{false}\FunctionTok{,}
    \DataTypeTok{"diffEditor.maxComputationTime"}\FunctionTok{:} \DecValTok{0}\FunctionTok{,}
    \DataTypeTok{"diffEditor.codeLens"}\FunctionTok{:} \KeywordTok{true}\FunctionTok{,}
    \DataTypeTok{"diffEditor.wordWrap"}\FunctionTok{:} \StringTok{"off"}\FunctionTok{,}
    
    \ErrorTok{//} \ErrorTok{==========================================}
    \ErrorTok{//} \ErrorTok{LANGUAGE{-}SPECIFIC}
    \ErrorTok{//} \ErrorTok{==========================================}
    \DataTypeTok{"[python]"}\FunctionTok{:} \FunctionTok{\{}
        \DataTypeTok{"editor.defaultFormatter"}\FunctionTok{:} \StringTok{"charliermarsh.ruff"}\FunctionTok{,}
        \DataTypeTok{"diffEditor.ignoreTrimWhitespace"}\FunctionTok{:} \KeywordTok{false}
    \FunctionTok{\},}
    \DataTypeTok{"[quarto]"}\FunctionTok{:} \FunctionTok{\{}
        \DataTypeTok{"editor.defaultFormatter"}\FunctionTok{:} \StringTok{"quarto.quarto"}
    \FunctionTok{\},}
    \DataTypeTok{"[markdown]"}\FunctionTok{:} \FunctionTok{\{}
        \DataTypeTok{"editor.defaultFormatter"}\FunctionTok{:} \StringTok{"esbenp.prettier{-}vscode"}
    \FunctionTok{\},}
    \DataTypeTok{"[yaml]"}\FunctionTok{:} \FunctionTok{\{}
        \DataTypeTok{"editor.defaultFormatter"}\FunctionTok{:} \StringTok{"quarto.quarto"}
    \FunctionTok{\},}
    \DataTypeTok{"[latex]"}\FunctionTok{:} \FunctionTok{\{}
        \DataTypeTok{"editor.defaultFormatter"}\FunctionTok{:} \StringTok{"James{-}Yu.latex{-}workshop"}
    \FunctionTok{\},}
    
    \ErrorTok{//} \ErrorTok{==========================================}
    \ErrorTok{//} \ErrorTok{QUARTO}
    \ErrorTok{//} \ErrorTok{==========================================}
    \DataTypeTok{"quarto.visualEditor.spellingDictionary"}\FunctionTok{:} \StringTok{"es\_ES"}\FunctionTok{,}
    
    \ErrorTok{//} \ErrorTok{==========================================}
    \ErrorTok{//} \ErrorTok{NOTEBOOKS}
    \ErrorTok{//} \ErrorTok{==========================================}
    \DataTypeTok{"notebook.output.wordWrap"}\FunctionTok{:} \KeywordTok{true}\FunctionTok{,}
    \DataTypeTok{"jupyter.interactiveWindow.textEditor.executeSelection"}\FunctionTok{:} \KeywordTok{true}\FunctionTok{,}
    
    \ErrorTok{//} \ErrorTok{==========================================}
    \ErrorTok{//} \ErrorTok{MISC}
    \ErrorTok{//} \ErrorTok{==========================================}
    \DataTypeTok{"debug.console.wordWrap"}\FunctionTok{:} \KeywordTok{false}\FunctionTok{,}
    \DataTypeTok{"chat.editor.wordWrap"}\FunctionTok{:} \StringTok{"on"}\FunctionTok{,}
    \DataTypeTok{"yaml.schemas"}\FunctionTok{:} \FunctionTok{\{}
        \DataTypeTok{"https://quarto.org/schemas/quarto.yaml"}\FunctionTok{:} \StringTok{"*.qmd"}
    \FunctionTok{\},}
    
    \ErrorTok{//} \ErrorTok{==========================================}
    \ErrorTok{//} \ErrorTok{POSITRON{-}SPECIFIC}
    \ErrorTok{//} \ErrorTok{==========================================}
    \DataTypeTok{"positron.assistant.enable"}\FunctionTok{:} \KeywordTok{true}\FunctionTok{,}
    \DataTypeTok{"positron.assistant.alwaysEnableApplyInEditorAction"}\FunctionTok{:} \KeywordTok{true}\FunctionTok{,}
    \DataTypeTok{"positron.assistant.alwaysIncludeCopilotTools"}\FunctionTok{:} \KeywordTok{true}\FunctionTok{,}
    \DataTypeTok{"positron.assistant.showTokenUsage.enable"}\FunctionTok{:} \KeywordTok{true}\FunctionTok{,}
    
    \ErrorTok{//} \ErrorTok{==========================================}
    \ErrorTok{//} \ErrorTok{ACCESSIBILITY} \ErrorTok{(ocultar} \ErrorTok{para} \ErrorTok{minimalista)}
    \ErrorTok{//} \ErrorTok{==========================================}
    \DataTypeTok{"accessibility.verbosity.terminal"}\FunctionTok{:} \KeywordTok{false}\FunctionTok{,}
    \DataTypeTok{"accessibility.hideAccessibleView"}\FunctionTok{:} \KeywordTok{true}\FunctionTok{,}
    
    \ErrorTok{//} \ErrorTok{==========================================}
    \ErrorTok{//} \ErrorTok{EXTENSIONES} \ErrorTok{(si} \ErrorTok{las} \ErrorTok{instalas)}
    \ErrorTok{//} \ErrorTok{==========================================}
    \ErrorTok{//} \ErrorTok{GitLens} \ErrorTok{(si} \ErrorTok{disponible)}
    \DataTypeTok{"gitlens.graph.minimap.additionalTypes"}\FunctionTok{:} \OtherTok{[}
        \StringTok{"localBranches"}\OtherTok{,}
        \StringTok{"stashes"}\OtherTok{,}
        \StringTok{"remoteBranches"}\OtherTok{,}
        \StringTok{"pullRequests"}\OtherTok{,}
        \StringTok{"tags"}
    \OtherTok{]}\FunctionTok{,}
    
    \ErrorTok{//} \ErrorTok{CodeSnap} \ErrorTok{(si} \ErrorTok{disponible)}
    \DataTypeTok{"codesnap.transparentBackground"}\FunctionTok{:} \KeywordTok{true}\FunctionTok{,}
    \DataTypeTok{"codesnap.showWindowTitle"}\FunctionTok{:} \KeywordTok{true}\FunctionTok{,}
    \DataTypeTok{"codesnap.containerPadding"}\FunctionTok{:} \StringTok{"0em"}\FunctionTok{,}
    \DataTypeTok{"codesnap.roundedCorners"}\FunctionTok{:} \KeywordTok{false}\FunctionTok{,}
    
    \ErrorTok{//} \ErrorTok{LaTeX} \ErrorTok{Workshop} \ErrorTok{(si} \ErrorTok{usas} \ErrorTok{LaTeX)}
    \DataTypeTok{"latex{-}workshop.formatting.latex"}\FunctionTok{:} \StringTok{"latexindent"}\FunctionTok{,}
    \DataTypeTok{"latex{-}workshop.latex.option.maxPrintLine.enabled"}\FunctionTok{:} \KeywordTok{false}\FunctionTok{,}
    
    \ErrorTok{//} \ErrorTok{Indenticator} \ErrorTok{(si} \ErrorTok{lo} \ErrorTok{instalas)}
    \DataTypeTok{"indenticator.width"}\FunctionTok{:} \FloatTok{0.5}\FunctionTok{,}
    \DataTypeTok{"indenticator.color.dark"}\FunctionTok{:} \StringTok{"rgba(255,255,255,0.1)"}
\FunctionTok{\}}
\end{Highlighting}
\end{Shaded}

\subsection{Extensiones a Instalar (desde Open
VSX)}\label{extensiones-a-instalar-desde-open-vsx}

Para replicar el setup de VS Code, se instala:

\begin{enumerate}
\def\labelenumi{\arabic{enumi}.}
\tightlist
\item
  \textbf{Material Icon Theme}
\item
  \textbf{One Dark Pro} (tema)
\item
  \textbf{Prettier} (formateo Markdown)
\item
  \textbf{LaTeX Workshop} (si usas LaTeX)
\item
  \textbf{GitLens} (si disponible)
\item
  \textbf{Better Comments} (opcional)
\item
  \textbf{Code Spell Checker} (opcional, español)
\end{enumerate}

\subsection{Notas sobre Extensiones No
Disponibles}\label{notas-sobre-extensiones-no-disponibles}

\textbf{No disponibles/innecesarias en Positron:}

\begin{itemize}
\tightlist
\item
  \texttt{animations} (vscode-animations): No compatible
\item
  \texttt{glassit}: No compatible (transparencia)
\item
  \texttt{apc.imports}: No compatible
\item
  Python extension: Ya incluida nativamente
\item
  GitHub Copilot: Usa Positron Assistant en su lugar
\end{itemize}

\section{Trucos y Tips}\label{trucos-y-tips}

\subsection{Command Palette}\label{command-palette}

\begin{verbatim}
Ctrl+Shift+P
\end{verbatim}

Todo está en el Command Palette. No recuerdas un atajo? Búscalo aquí.

\subsection{Quick Open}\label{quick-open}

\begin{verbatim}
Ctrl+P
\end{verbatim}

Abre archivos rápidamente sin usar el Explorer.

\subsection{Multi-Cursor}\label{multi-cursor}

\begin{verbatim}
Alt+Click: Añadir cursor
Ctrl+Alt+Up/Down: Añadir cursor arriba/abajo
Ctrl+D: Seleccionar siguiente ocurrencia
Ctrl+Shift+L: Seleccionar todas las ocurrencias
\end{verbatim}

\subsection{Snippets}\label{snippets}

Crea tus propios snippets en:

\begin{verbatim}
File → Preferences → Configure User Snippets
\end{verbatim}

\textbf{Ejemplo Python:}

\begin{Shaded}
\begin{Highlighting}[]
\FunctionTok{\{}
    \DataTypeTok{"Import pandas"}\FunctionTok{:} \FunctionTok{\{}
        \DataTypeTok{"prefix"}\FunctionTok{:} \StringTok{"imp"}\FunctionTok{,}
        \DataTypeTok{"body"}\FunctionTok{:} \OtherTok{[}
            \StringTok{"import pandas as pd"}\OtherTok{,}
            \StringTok{"import numpy as np"}\OtherTok{,}
            \StringTok{"import matplotlib.pyplot as plt"}
        \OtherTok{]}
    \FunctionTok{\}}
\FunctionTok{\}}
\end{Highlighting}
\end{Shaded}

\subsection{Integrated Terminal
Profiles}\label{integrated-terminal-profiles}

Configura perfiles para diferentes shells o entornos virtuales.

\subsection{Workspace Settings}\label{workspace-settings}

Para proyectos específicos, usa workspace settings:

\begin{verbatim}
.vscode/settings.json
\end{verbatim}

\subsection{Fold/Unfold Code}\label{foldunfold-code}

\begin{verbatim}
Ctrl+Shift+[: Fold
Ctrl+Shift+]: Unfold
Ctrl+K Ctrl+0: Fold all
Ctrl+K Ctrl+J: Unfold all
\end{verbatim}

\subsection{Breadcrumbs}\label{breadcrumbs}

Si los necesitas:

\begin{Shaded}
\begin{Highlighting}[]
\FunctionTok{\{}
    \DataTypeTok{"breadcrumbs.enabled"}\FunctionTok{:} \KeywordTok{true}
\FunctionTok{\}}
\end{Highlighting}
\end{Shaded}

\subsection{Integrated Git Diff}\label{integrated-git-diff}

\begin{verbatim}
Click en archivo modificado en Source Control
\end{verbatim}

\subsection{Sync Settings}\label{sync-settings}

Positron puede sincronizar settings entre máquinas (similar a VS Code
Settings Sync).

\subsection{Trucos prácticos que uso día a
día}\label{trucos-pruxe1cticos-que-uso-duxeda-a-duxeda}

\textbf{1. Ejecutar código selectivamente sin salir del editor}

\begin{itemize}
\tightlist
\item
  \texttt{Ctrl\ +\ Enter} → Ejecuta la línea actual o selección
\item
  \texttt{Ctrl\ +\ Alt\ +\ Home} → Ejecuta \textbf{todo desde el inicio
  del archivo hasta la línea actual} (ideal para setups y carga de
  datos)
\item
  \texttt{Ctrl\ +\ Alt\ +\ End} → Ejecuta \textbf{desde la línea actual
  hasta el final} (muy útil para pruebas rápidas al final del script)
\item
  \texttt{Ctrl\ +\ Shift\ +\ P} (o \texttt{Cmd} en macOS) → Ejecuta
  \textbf{archivo completo} en la consola
\end{itemize}

\textbf{Truco pro}: Combina \texttt{Ctrl\ +\ Alt\ +\ Home} +
\texttt{Ctrl\ +\ Enter} para rehacer rápidamente todo el análisis desde
cero sin tocar el ratón.

\textbf{2. Navegación ultrarrápida entre editor y consola}

\begin{Shaded}
\begin{Highlighting}[]
\NormalTok{Ctrl + K → F     → Enfocar consola (muy rápido para escribir rápido)}
\NormalTok{Ctrl + K → V     → Enfocar panel de Variables}
\NormalTok{Ctrl + K → R     → Enfocar panel de Help (documentación)}
\NormalTok{Ctrl + L         → Limpiar la consola (equivalente al clásico Ctrl + L de RStudio)}
\NormalTok{Ctrl + Shift + 0 → Reiniciar el intérprete (cuando todo se pone loco)}
\end{Highlighting}
\end{Shaded}

\textbf{Truco diario}: Crea un flujo mental ``editor → consola →
variables'' y usa estos tres atajos en secuencia rapidísima.

\textbf{3. El workflow más rápido para explorar dataframes (mi
favorito)}

\begin{enumerate}
\def\labelenumi{\arabic{enumi}.}
\item
  Cargas datos en consola o script:

\begin{Shaded}
\begin{Highlighting}[]
\ImportTok{import}\NormalTok{ pandas }\ImportTok{as}\NormalTok{ pd}
\NormalTok{df }\OperatorTok{=}\NormalTok{ pd.read\_csv(}\StringTok{"datos\_grandes.csv"}\NormalTok{)}
\end{Highlighting}
\end{Shaded}

  o

\begin{Shaded}
\begin{Highlighting}[]
\NormalTok{df }\OtherTok{\textless{}{-}} \FunctionTok{read\_csv}\NormalTok{(}\StringTok{"datos\_grandes.csv"}\NormalTok{)}
\end{Highlighting}
\end{Shaded}
\item
  Opciones rápidas de exploración:

  \begin{itemize}
  \item
    Doble clic en \texttt{df} en el \textbf{Variables Pane} → abre
    \textbf{Data Explorer} completo (filtrar, ordenar, buscar,
    estadísticas)
  \item
    \texttt{View(df)} (R) o \texttt{df.head(20)} (Python) → vista rápida
    en consola
  \item
    Clic derecho en variable → ``View in Data Explorer''
  \item
    Atajo personalizado recomendado (agrega a keybindings.json):

\begin{Shaded}
\begin{Highlighting}[]
\FunctionTok{\{}
    \DataTypeTok{"key"}\FunctionTok{:} \StringTok{"f2"}\FunctionTok{,}
    \DataTypeTok{"command"}\FunctionTok{:} \StringTok{"positron.dataExplorer.openDataExplorer"}\FunctionTok{,}
    \DataTypeTok{"when"}\FunctionTok{:} \StringTok{"positronVariablesFocus"}
\FunctionTok{\}}
\end{Highlighting}
\end{Shaded}
  \end{itemize}
\end{enumerate}

\textbf{4. Trucos rápidos de transformación de texto}

Positron hereda estas joyas de Code OSS:

\begin{itemize}
\tightlist
\item
  \textbf{Mayúsculas / minúsculas / título / oración}:

  \begin{enumerate}
  \def\labelenumi{\arabic{enumi}.}
  \tightlist
  \item
    Selecciona texto
  \item
    \texttt{Ctrl\ +\ Shift\ +\ P} → escribe ``Transform to''
  \item
    Elige: Upper Case / Lower Case / Title Case / Sentence Case
  \end{enumerate}
\item
  \textbf{Multi-cursor express} (imprescindible para renombrar
  columnas):

  \begin{itemize}
  \tightlist
  \item
    \texttt{Ctrl\ +\ D} → selecciona siguiente ocurrencia
  \item
    \texttt{Ctrl\ +\ Shift\ +\ L} → selecciona \textbf{todas} las
    ocurrencias en el archivo
  \item
    \texttt{Alt\ +\ Click} → múltiples cursores manuales
  \end{itemize}
\item
  \textbf{Truco ultra-rápido}: \texttt{Ctrl\ +\ Shift\ +\ P} → ``Join
  Lines'' (une líneas seleccionadas en una) → genial para convertir
  listas largas en una sola línea de código.
\end{itemize}

\textbf{5. Cambio rápido entre Python y R en el mismo proyecto}

\begin{itemize}
\tightlist
\item
  Usa el selector de intérprete en la \textbf{barra superior} (junto al
  nombre del archivo)
\item
  Puedes tener \textbf{dos sesiones simultáneas}: una Python y otra R
\item
  Truco práctico:

  \begin{enumerate}
  \def\labelenumi{\arabic{enumi}.}
  \tightlist
  \item
    Cargas datos en Python (pandas)
  \item
    Cambias a sesión R
  \item
    Usas \texttt{reticulate::py\$df} para traer el dataframe de Python a
    R
  \item
    Sigues limpiando con tidyverse
  \end{enumerate}
\end{itemize}

\textbf{6. Trucos con Plots Pane que ahorran mucho tiempo}

\begin{itemize}
\tightlist
\item
  Historial automático de gráficos → flechas arriba/abajo para navegar
\item
  Clic derecho en plot → Exportar (PNG/SVG/PDF) rápido
\item
  Truco: \texttt{plt.show()} (Python) o \texttt{print()} (ggplot2) →
  aparece automáticamente
\item
  Usa \texttt{plt.ion()} (modo interactivo matplotlib) para actualizar
  gráficos en vivo sin repetir \texttt{show()}
\end{itemize}

\textbf{7. Bonus: Mini-flujo de análisis exprés}

\begin{Shaded}
\begin{Highlighting}[]
\NormalTok{1. Ctrl + N → nuevo script}
\NormalTok{2. Pega/pega datos o carga rápida}
\NormalTok{3. Ctrl + Alt + Home → ejecuta todo hasta cursor}
\NormalTok{4. Doble clic variable → Data Explorer}
\NormalTok{5. F1 sobre función → Help pane}
\NormalTok{6. Selecciona gráfico → Plots pane}
\NormalTok{7. Ctrl + Shift + K → renderiza Quarto (si estás en .qmd)}
\end{Highlighting}
\end{Shaded}

\section{Solución de Problemas}\label{soluciuxf3n-de-problemas}

\subsection{Problema: Python interpreter no
detectado}\label{problema-python-interpreter-no-detectado}

\textbf{Solución:}

\begin{verbatim}
Command Palette → Python: Select Interpreter
\end{verbatim}

O especificar manualmente:

\begin{Shaded}
\begin{Highlighting}[]
\FunctionTok{\{}
    \DataTypeTok{"python.defaultInterpreterPath"}\FunctionTok{:} \StringTok{"/path/to/python"}
\FunctionTok{\}}
\end{Highlighting}
\end{Shaded}

\subsection{Problema: R no funciona}\label{problema-r-no-funciona}

\textbf{Verificar instalación:}

\begin{Shaded}
\begin{Highlighting}[]
\FunctionTok{which}\NormalTok{ R}
\ExtensionTok{R} \AttributeTok{{-}{-}version}
\end{Highlighting}
\end{Shaded}

\textbf{Instalar paquetes necesarios:}

\begin{Shaded}
\begin{Highlighting}[]
\NormalTok{pak}\SpecialCharTok{::}\FunctionTok{pak}\NormalTok{(}\FunctionTok{c}\NormalTok{(}\StringTok{"usethis"}\NormalTok{, }\StringTok{"cli"}\NormalTok{, }\StringTok{"crayon"}\NormalTok{, }\StringTok{"rlang"}\NormalTok{))}
\end{Highlighting}
\end{Shaded}

\subsection{Problema: Console no
inicia}\label{problema-console-no-inicia}

\textbf{Reiniciar intérprete:}

\begin{verbatim}
Ctrl+Shift+0
\end{verbatim}

\textbf{Verificar IPython:}

\begin{Shaded}
\begin{Highlighting}[]
\ExtensionTok{pip}\NormalTok{ install ipython}
\end{Highlighting}
\end{Shaded}

\subsection{Problema: Extensiones no
aparecen}\label{problema-extensiones-no-aparecen}

\textbf{Recordar:} Positron usa Open VSX, no VS Marketplace. Algunas
extensiones pueden no estar disponibles.

\textbf{Alternativa:} Instalar .vsix manualmente desde GitHub de la
extensión.

\subsection{Problema: Terminal no hereda
PATH}\label{problema-terminal-no-hereda-path}

\begin{Shaded}
\begin{Highlighting}[]
\FunctionTok{\{}
    \DataTypeTok{"terminal.integrated.inheritEnv"}\FunctionTok{:} \KeywordTok{false}
\FunctionTok{\}}
\end{Highlighting}
\end{Shaded}

O agregar al shell config (\textasciitilde/.bashrc,
\textasciitilde/.zshrc, etc.):

\begin{Shaded}
\begin{Highlighting}[]
\BuiltInTok{export} \VariableTok{PATH}\OperatorTok{=}\StringTok{"/path/to/anaconda3/bin:}\VariableTok{$PATH}\StringTok{"}
\end{Highlighting}
\end{Shaded}

\subsection{Problema: Plots no
aparecen}\label{problema-plots-no-aparecen}

\textbf{Python:}

\begin{Shaded}
\begin{Highlighting}[]
\ImportTok{import}\NormalTok{ matplotlib.pyplot }\ImportTok{as}\NormalTok{ plt}
\NormalTok{plt.ion()  }\CommentTok{\# Interactive mode}
\end{Highlighting}
\end{Shaded}

\textbf{R:}

\begin{Shaded}
\begin{Highlighting}[]
\FunctionTok{options}\NormalTok{(}\AttributeTok{device =} \StringTok{"quartz"}\NormalTok{)  }\CommentTok{\# macOS}
\CommentTok{\# o}
\FunctionTok{options}\NormalTok{(}\AttributeTok{device =} \StringTok{"x11"}\NormalTok{)  }\CommentTok{\# Linux}
\end{Highlighting}
\end{Shaded}

\subsection{Problema: Quarto no
renderiza}\label{problema-quarto-no-renderiza}

\textbf{Verificar instalación:}

\begin{Shaded}
\begin{Highlighting}[]
\ExtensionTok{quarto}\NormalTok{ check}
\end{Highlighting}
\end{Shaded}

\textbf{Reinstalar extensión:}

\begin{verbatim}
Extensions → Quarto → Reload
\end{verbatim}

\subsection{Problema: Settings no se
guardan}\label{problema-settings-no-se-guardan}

\textbf{Ubicación correcta:}

\begin{verbatim}
~/.config/Positron/User/settings.json
\end{verbatim}

\textbf{Permisos:}

\begin{Shaded}
\begin{Highlighting}[]
\FunctionTok{chmod}\NormalTok{ 644 \textasciitilde{}/.config/Positron/User/settings.json}
\end{Highlighting}
\end{Shaded}

\subsection{Problema: Performance
lenta}\label{problema-performance-lenta}

\textbf{Deshabilitar extensiones innecesarias}

\textbf{Reducir history:}

\begin{Shaded}
\begin{Highlighting}[]
\FunctionTok{\{}
    \DataTypeTok{"terminal.integrated.scrollback"}\FunctionTok{:} \DecValTok{1000}
\FunctionTok{\}}
\end{Highlighting}
\end{Shaded}

\textbf{Deshabilitar features:}

\begin{Shaded}
\begin{Highlighting}[]
\FunctionTok{\{}
    \DataTypeTok{"editor.minimap.enabled"}\FunctionTok{:} \KeywordTok{false}\FunctionTok{,}
    \DataTypeTok{"editor.suggest.preview"}\FunctionTok{:} \KeywordTok{false}
\FunctionTok{\}}
\end{Highlighting}
\end{Shaded}

\section{Conclusión}\label{conclusiuxf3n}

\subsection{Resumen de Cambios de VS Code a
Positron}\label{resumen-de-cambios-de-vs-code-a-positron}

\textbf{Lo Mismo:}

\begin{itemize}
\tightlist
\item
  Editor de código
\item
  Atajos de teclado (95\%)
\item
  Terminal
\item
  Extensions system
\item
  Git integration
\item
  Settings system
\end{itemize}

\textbf{Lo Nuevo:}

\begin{itemize}
\tightlist
\item
  Console (IPython/IRKernel)
\item
  Variables Pane
\item
  Data Explorer
\item
  Plots Pane
\item
  Help Pane
\item
  Connections Pane
\item
  Soporte Python/R nativo
\item
  AI Assistant (Positron Assistant, Databot)
\end{itemize}

\textbf{Lo Diferente:}

\begin{itemize}
\tightlist
\item
  Open VSX en lugar de VS Marketplace
\item
  Algunas extensiones no disponibles
\item
  Mejor integración para data science
\end{itemize}

\subsection{Próximos Pasos}\label{pruxf3ximos-pasos}

\begin{enumerate}
\def\labelenumi{\arabic{enumi}.}
\tightlist
\item
  \textbf{Instalar Positron}
\item
  \textbf{Importar settings de VS Code}
\item
  \textbf{Configurar Python/R}
\item
  \textbf{Instalar extensiones esenciales}
\item
  \textbf{Personalizar tema y layout}
\item
  \textbf{Explorar Console y panes}
\item
  \textbf{Crear tu primer documento Quarto}
\end{enumerate}

\subsection{Recursos}\label{recursos}

\textbf{Documentación:}

\begin{itemize}
\tightlist
\item
  https://positron.posit.co/
\item
  https://github.com/posit-dev/positron
\end{itemize}

\textbf{Comunidad:}

\begin{itemize}
\tightlist
\item
  GitHub Discussions: https://github.com/posit-dev/positron/discussions
\item
  Posit Community: https://community.posit.co/
\end{itemize}

\textbf{Aprendizaje:}

\begin{itemize}
\tightlist
\item
  Quarto: https://quarto.org/
\item
  Python: https://docs.python.org/
\item
  R: https://www.r-project.org/
\end{itemize}

\section{Publicaciones Similares}\label{publicaciones-similares}

Si te interesó este artículo, te recomendamos que explores otros blogs y
recursos relacionados que pueden ampliar tus conocimientos. Aquí te dejo
algunas sugerencias:

\begin{enumerate}
\def\labelenumi{\arabic{enumi}.}
\tightlist
\item
  \href{https://chaska-x.netlify.app/operating-system/2017-05-21-comandos-de-informacion-windows/index.pdf}{\faIcon{file-pdf}}
  \href{https://chaska-x.netlify.app/operating-system/2017-05-21-comandos-de-informacion-windows}{Comandos
  De Informacion Windows}
\item
  \href{https://chaska-x.netlify.app/operating-system/2019-06-19-adb/index.pdf}{\faIcon{file-pdf}}
  \href{https://chaska-x.netlify.app/operating-system/2019-06-19-adb}{Adb}
\item
  \href{https://chaska-x.netlify.app/operating-system/2021-08-17-limpieza-y-optimizacion-de-pc/index.pdf}{\faIcon{file-pdf}}
  \href{https://chaska-x.netlify.app/operating-system/2021-08-17-limpieza-y-optimizacion-de-pc}{Limpieza
  Y Optimizacion De Pc}
\item
  \href{https://chaska-x.netlify.app/operating-system/2021-10-21-usando-apk-en-windown-11/index.pdf}{\faIcon{file-pdf}}
  \href{https://chaska-x.netlify.app/operating-system/2021-10-21-usando-apk-en-windown-11}{Usando
  Apk En Windown 11}
\item
  \href{https://chaska-x.netlify.app/operating-system/2022-05-12-gestionar-versiones-de-jdk-en-kubuntu/index.pdf}{\faIcon{file-pdf}}
  \href{https://chaska-x.netlify.app/operating-system/2022-05-12-gestionar-versiones-de-jdk-en-kubuntu}{Gestionar
  Versiones De Jdk En Kubuntu}
\item
  \href{https://chaska-x.netlify.app/operating-system/2022-07-21-instalar-tor-browser/index.pdf}{\faIcon{file-pdf}}
  \href{https://chaska-x.netlify.app/operating-system/2022-07-21-instalar-tor-browser}{Instalar
  Tor Browser}
\item
  \href{https://chaska-x.netlify.app/operating-system/2022-08-14-crear-enlaces-duros-o-hard-link-en-linux/index.pdf}{\faIcon{file-pdf}}
  \href{https://chaska-x.netlify.app/operating-system/2022-08-14-crear-enlaces-duros-o-hard-link-en-linux}{Crear
  Enlaces Duros O Hard Link En Linux}
\item
  \href{https://chaska-x.netlify.app/operating-system/2022-09-27-comandos-vim/index.pdf}{\faIcon{file-pdf}}
  \href{https://chaska-x.netlify.app/operating-system/2022-09-27-comandos-vim}{Comandos
  Vim}
\item
  \href{https://chaska-x.netlify.app/operating-system/2023-02-16-guia-de-git-y-github/index.pdf}{\faIcon{file-pdf}}
  \href{https://chaska-x.netlify.app/operating-system/2023-02-16-guia-de-git-y-github}{Guia
  De Git Y Github}
\item
  \href{https://chaska-x.netlify.app/operating-system/2023-05-02-00-primeros-pasos-en-linux/index.pdf}{\faIcon{file-pdf}}
  \href{https://chaska-x.netlify.app/operating-system/2023-05-02-00-primeros-pasos-en-linux}{00
  Primeros Pasos En Linux}
\item
  \href{https://chaska-x.netlify.app/operating-system/2023-06-17-01-introduccion-linux/index.pdf}{\faIcon{file-pdf}}
  \href{https://chaska-x.netlify.app/operating-system/2023-06-17-01-introduccion-linux}{01
  Introduccion Linux}
\item
  \href{https://chaska-x.netlify.app/operating-system/2023-06-18-02-distribuciones-linux/index.pdf}{\faIcon{file-pdf}}
  \href{https://chaska-x.netlify.app/operating-system/2023-06-18-02-distribuciones-linux}{02
  Distribuciones Linux}
\item
  \href{https://chaska-x.netlify.app/operating-system/2023-06-19-03-instalacion-linux/index.pdf}{\faIcon{file-pdf}}
  \href{https://chaska-x.netlify.app/operating-system/2023-06-19-03-instalacion-linux}{03
  Instalacion Linux}
\item
  \href{https://chaska-x.netlify.app/operating-system/2023-06-20-04-administracion-particiones-volumenes/index.pdf}{\faIcon{file-pdf}}
  \href{https://chaska-x.netlify.app/operating-system/2023-06-20-04-administracion-particiones-volumenes}{04
  Administracion Particiones Volumenes}
\item
  \href{https://chaska-x.netlify.app/operating-system/2023-07-01-atajos-de-teclado-y-comandos-para-usar-vim/index.pdf}{\faIcon{file-pdf}}
  \href{https://chaska-x.netlify.app/operating-system/2023-07-01-atajos-de-teclado-y-comandos-para-usar-vim}{Atajos
  De Teclado Y Comandos Para Usar Vim}
\item
  \href{https://chaska-x.netlify.app/operating-system/2024-07-15-instalando-specitify/index.pdf}{\faIcon{file-pdf}}
  \href{https://chaska-x.netlify.app/operating-system/2024-07-15-instalando-specitify}{Instalando
  Specitify}
\item
  \href{https://chaska-x.netlify.app/operating-system/2025-07-10-gestiona-tus-dotfiles-con-gnu-stow/index.pdf}{\faIcon{file-pdf}}
  \href{https://chaska-x.netlify.app/operating-system/2025-07-10-gestiona-tus-dotfiles-con-gnu-stow}{Gestiona
  Tus Dotfiles Con Gnu Stow}
\end{enumerate}

Esperamos que encuentres estas publicaciones igualmente interesantes y
útiles. ¡Disfruta de la lectura!






\end{document}
